\chapter{Sectional Curvature Comparison I}
\label{chap:pet-sectional-curvature-comparison-i}

The chapter develops covariant differentiation along curves, Jacobi fields, and
Synge's second variation formula, then applies Riccati and Rauch comparison to
nonpositive and positive sectional curvature.

\section{The Connection Along Curves}

\begin{lemma}[the chart--connection bridge]
  \label{lem:chart-connection-bridge}
  \lean{PetersenLib.leviCivita_chartFrame_christoffel, PetersenLib.exists_chartFrame_leviCivita_christoffel_nhds}
  \leanok
  Infrastructure, not a numbered result of the source; recorded because the
  formalization needs it and Petersen takes it for granted.

  Around any $p$ in a chart around $\alpha$ there is a local frame $Z_1,\dots,Z_n$
  and a neighbourhood $U\ni p$ on which the Levi-Civita connection of $g$ is
  computed by the Christoffel symbols of the second kind:
  \[
    \nabla_{Z_i}Z_j=\sum_m\Gamma^m_{ij}\,Z_m,\qquad
    \Gamma^m_{ij}=\tfrac12\sum_lG^{ml}\big(\partial_iG_{lj}+\partial_jG_{li}
      -\partial_lG_{ij}\big).
  \]
  Here $\nabla$ is the abstract (Koszul) connection through which the curvature
  tensor $R$ and $\sec$ are defined, while $\Gamma$ is the chart symbol through
  which geodesics, $\exp$, and the variational calculus are defined. Petersen
  moves between the two without comment; in the formalization they are distinct
  objects, and every result of this chapter that couples curvature to geodesics
  must identify them.

  \emph{Remark on how this is actually used.} The identification that
  \cref{lem:pet-ch6-third-partial-commutator} consumes is the sharper
  \cref{lem:curvature-chart-bridge}, which reaches all the way to $R$ rather than
  stopping at $\nabla$, and which is obtained from
  \cref{prop:pet-ch3-curvature-coordinate-formula} rather than from this lemma.
  This lemma remains the general-purpose statement of the identification (and is
  strictly more general than its Chapter 3 counterpart, being stated for an
  arbitrary chart centre $\alpha$ rather than only for $\alpha=p$).

  The proof also needs germ-locality of $\nabla$ (that $\nabla_vX$ depends only on
  the germ of $X$), which the source likewise leaves implicit.
\end{lemma}

\begin{lemma}[the curvature--chart bridge]
  \label{lem:curvature-chart-bridge}
  \lean{PetersenLib.chartCurvatureContraction2_eq_neg_curvatureTensorAt}
  \uses{prop:pet-ch3-curvature-coordinate-formula}
  \leanok
  Infrastructure, not a numbered result of the source; recorded because the
  formalization needs it and Petersen takes it for granted.

  Let $R_{\mathrm{chart}}(X,Y)Z=\sum_l\big(\sum_{i,j,k}R^l{}_{ijk}(y)X^iY^jZ^k\big)e_l$
  be the curvature contraction written directly in the Christoffel symbols of the
  chart at $p$, with
  \[
    R^l{}_{ijk}=\partial_j\Gamma^l_{ik}-\partial_i\Gamma^l_{jk}
      +\sum_s\big(\Gamma^s_{ik}\Gamma^l_{js}-\Gamma^s_{jk}\Gamma^l_{is}\big).
  \]
  Then, at the centre of that chart,
  \[
    R_{\mathrm{chart}}(X,Y)Z=-R(X,Y)Z,
  \]
  with $R$ the abstract (Koszul) curvature tensor of \cref{def:pet-ch3-curvature-tensor}.

  Two remarks on the shape of the statement.

  \emph{The sign.} $R_{\mathrm{chart}}$ is written in do Carmo's convention
  $R(X,Y)=\nabla_Y\nabla_X-\nabla_X\nabla_Y+\nabla_{[X,Y]}$, the negative of the
  convention of \cref{def:pet-ch3-curvature-tensor}; the two coefficient
  expressions are exact negations with the same index slots. In
  \cref{lem:pet-ch6-third-partial-commutator} this sign cancels against the
  opposite parameter order in which the coordinate commutator is proved, so
  Petersen's identity holds with Petersen's $R$ and no stray sign.

  \emph{Why the chart centre suffices.}
  \cref{prop:pet-ch3-curvature-coordinate-formula} is stated on the diagonal —
  chart centre $=$ evaluation point. This is no restriction for
  \cref{lem:pet-ch6-third-partial-commutator}, whose conclusion is pointwise at
  $p=c(u_0,s_0,t_0)$: read the map in the chart at that very point. The step that
  makes the diagonal case collapse is that at the centre of its own chart the
  chart frame $\partial_i|_p$ \emph{is} the model basis $e_i$, so that the
  chart-frame expansion of a tangent vector at $p$ is the ordinary basis expansion
  of $E$ and the coordinate-to-abstract comparison map is the identity.
\end{lemma}

\begin{definition}[covariant derivative along a curve]
  \label{def:pet-ch6-vector-field-along-curve}
  \lean{PetersenLib.derivAlongCurve,PetersenLib.covariantDerivCoord_transfer}
  \source{petersen:page-0248,petersen:page-0249}
  \leanok
  For a curve $c:I\to M$, a \emph{vector field along $c$} is a map $V:I\to TM$
  with $V(t)\in T_{c(t)}M$ for all $t\in I$. Thinking of $V$ as the variational
  field of a variation $\bar c:(-\varepsilon,\varepsilon)\times I\to M$ with
  $\bar c(0,t)=c(t)$, $\frac{\partial\bar c}{\partial s}(0,t)=V(t)$, the
  \emph{covariant derivative of $V$ along $c$} is defined by
  $\dot V(t)=\frac{d}{dt}V(t)=\frac{\partial^2\bar c}{\partial t\,\partial s}(0,t)$.
  In local coordinates, writing $V(t)=V^k(t)\partial_k$ and
  $\bar c^k(s,t)=\tfrac{\partial\bar c^k}{\partial s}(0,t)$, one computes
  \[
    \frac{\partial^2\bar c}{\partial t\,\partial s}(0,t)
    =\dot V^k(t)\partial_k+V^i(t)\dot c^j(t)\Gamma^k_{ij}\partial_k,
  \]
  which is manifestly independent of the choice of variation and of the
  coordinate system, and agrees with $\nabla_{\dot c(t_0)}X$ whenever
  $V(t)=X_{c(t)}$ for a vector field $X$ near $c(t_0)$. In general $\dot V(t_0)$
  need not vanish when $\dot c(t_0)=0$ (e.g.\ for a constant curve $c$, where $V$
  is simply a curve in the fixed space $T_{c(t_0)}M$). Selecting a two-parameter
  variation as needed for a second field $W$ along $c$, the local formula also
  yields the product rule $\frac{d}{dt}g(V,W)=g(\dot V,W)+g(V,\dot W)$ and
  linearity $\frac{d}{dt}(V+W)=\dot V+\dot W$,
  $\frac{d}{dt}(\lambda V)=\dot\lambda V+\lambda\dot V$ for $\lambda:I\to\mathbb{R}$.
  When $M\subset\tilde M$, the derivative along $c$ in $M$ is obtained by
  computing $\dot V\in T\tilde M$ and orthogonally projecting onto $TM$.
\end{definition}

\begin{definition}[third and higher partial derivatives]
  \label{def:pet-ch6-higher-order-partials}
  \lean{PetersenLib.surfaceCovariantDeriv,PetersenLib.thirdPartialDerivative}
  \source{petersen:page-0249}
  \uses{def:pet-ch6-vector-field-along-curve}
  \leanok
  Using \cref{def:pet-ch6-vector-field-along-curve}, higher mixed partials of a
  multi-parameter map $c$ are defined recursively: to compute
  $\frac{\partial^3c}{\partial s\,\partial t\,\partial u}(s_0,t_0,u_0)$, set
  $V(s)=\frac{\partial^2c}{\partial t\,\partial u}(s,t_0,u_0)$ and define the
  third partial as $\frac{dV}{ds}(s_0)$. Because this recursive definition fixes
  a specific order of differentiation, there is no a priori reason for third
  partials to be symmetric in all three variables; it remains true, however,
  that $\frac{\partial^3c}{\partial s\,\partial t\,\partial u}
  =\frac{\partial^3c}{\partial s\,\partial u\,\partial t}$ (the last two
  variables, which enter only through the already-defined second mixed
  partial, may be swapped).
\end{definition}

\begin{example}[mixed partials on $S^2(1)$]
  \label{ex:pet-ch6-sphere-third-partials}
  \lean{PetersenLib.sphereThirdPartialsExample}
  \source{petersen:page-0250}
  \uses{def:pet-ch6-higher-order-partials}
  \dcref{ch6:6.1.1}
  \leanok
  Parametrize $S^2(1)\subset\mathbb{R}^3$ as a surface of revolution about the
  $x$-axis by $c(t,\theta)=(\cos t,\sin t\cos\theta,\sin t\sin\theta)$, and
  compute all derivatives in $\mathbb{R}^3$ before projecting onto $S^2(1)$. The
  curves $t\mapsto c(t,\theta)$ are geodesics, since
  $\partial_tc=(-\sin t,\cos t\cos\theta,\cos t\sin\theta)\in TS^2(1)$ while
  $\partial_t^2c=(-\cos t,-\sin t\cos\theta,-\sin t\sin\theta)=-c$ is normal to
  $S^2(1)$ and so projects to zero. Next,
  $\partial_\theta\partial_tc=(0,-\cos t\sin\theta,\cos t\cos\theta)$ is already
  tangent to $S^2(1)$, hence equals the intrinsic mixed partial. Finally
  \[
    \partial_t\partial_\theta\partial_tc=(0,\sin t\sin\theta,-\sin t\cos\theta)
    =\partial_\theta\partial_t^2c
  \]
  in $\mathbb{R}^3$, and both are tangent to $S^2(1)$; the first is
  $\partial_t\partial_\theta\partial_tc$ as computed intrinsically in $S^2(1)$,
  while the second has no intrinsic meaning as written, since $\partial_t^2c$
  must first be projected to $S^2(1)$ (where it vanishes) before differentiating
  in $\theta$. Hence intrinsically
  $\partial_\theta\partial_t^2c=0\ne\partial_t\partial_\theta\partial_tc$ along
  $t=\pi/2$, illustrating that third partials computed with different orderings of
  differentiation need not agree. Along the equator $t=\pi/2$,
  $\partial_\theta\partial_tc=0$.
\end{example}

\begin{lemma}[third partials and curvature]
  \label{lem:pet-ch6-third-partial-commutator}
  \lean{PetersenLib.thirdPartialCurvatureFormula}
  \source{petersen:page-0251}
  \uses{def:pet-ch6-higher-order-partials, lem:curvature-chart-bridge}
  \leanok
  \dcref{ch6:6.1.2}
  For a map $c(u,s,t)$ with values in $M$,
  \[
    \frac{\partial^3c}{\partial u\,\partial s\,\partial t}
    -\frac{\partial^3c}{\partial s\,\partial u\,\partial t}
    =R\!\left(\frac{\partial c}{\partial u},\frac{\partial c}{\partial s}\right)
    \frac{\partial c}{\partial t}.
  \]

  \emph{Formalization note.} The Lean statement is the identity for the
  $(u,s)$-surface with an \emph{arbitrary} $C^2$ field $V$ along it, of which
  Petersen's statement is the case $V=\frac{\partial c}{\partial t}$: the $t$-slot
  enters only through $V$, and the general form is what the Jacobi equation
  (\cref{def:pet-ch6-jacobi-field}) actually consumes. The surface is read in the
  chart centred at the base point $p=c(u_0,s_0,t_0)$, which by
  \cref{lem:curvature-chart-bridge} costs no generality, and the regularity Petersen
  leaves implicit is explicit.
\end{lemma}
\begin{proof}
  \leanok
  \sketch
  Use normal coordinates. Expanding the recursively defined third derivatives cancels ordinary coordinate third partials; the remaining Christoffel derivatives combine, via the curvature--chart bridge, to \(R(\partial_u c,\partial_s c)\partial_t c\). Keeping the quadratic terms instead proves the arbitrary-field form used by Lean.
\end{proof}

\begin{definition}[parallel field along a curve]
  \label{def:pet-ch6-parallel-field}
  \lean{PetersenLib.IsParallelAlong}
  \source{petersen:page-0252}
  \uses{def:pet-ch6-vector-field-along-curve}
  \leanok
  A vector field $V$ along $c$ is \emph{parallel along $c$} if $\dot V\equiv0$.
  The tangent field $\dot c$ along a geodesic is parallel, and (by
  \cref{ex:pet-ch6-sphere-third-partials}) the unit field perpendicular to a
  great circle in $S^2(1)$ is parallel. If $V,W$ are parallel along $c$ then
  $g(V,W)$ is constant along $c$ by the product rule, so parallel fields neither
  change length nor relative angle, generalizing parallel translation in
  Euclidean space.
\end{definition}

\begin{theorem}[existence and uniqueness of parallel fields]
  \label{thm:pet-ch6-parallel-field-existence-uniqueness}
  \lean{PetersenLib.parallelField_existence_uniqueness_global, PetersenLib.exists_isParallelSolOn_Ioo_of_openHyp}
  \source{petersen:page-0252}
  \uses{def:pet-ch6-parallel-field, thm:pet-ch6-parallel-field-chart-local}
  \leanok
  \dcref{ch6:6.1.3a}
  If $t_0\in I$ and $v\in T_{c(t_0)}M$, there is a unique parallel field $V(t)$
  defined on all of $I$ with $V(t_0)=v$.

  Global existence on the whole of $I$ — not merely locally — is the substance of
  the theorem: it is what distinguishes the \emph{linear} parallel-transport
  equation from the nonlinear geodesic equation, whose solutions exist only
  locally. The formalization assumes only that $c$ is $C^2$ on $I$; no
  chart-containment hypothesis is imposed, and
  \texttt{exists\_isParallelSolOn\_Ioo\_of\_openHyp} takes Petersen's hypothesis
  verbatim, with $c$ regular only on the \emph{open} interval.
\end{theorem}
\begin{proof}
  \uses{thm:pet-ch6-parallel-field-chart-local,def:pet-ch6-vector-field-along-curve}
  \leanok
  \sketch
  In a chart, parallel transport is the linear system \(\dot V^j=-\sum_i\alpha_i^jV^i\), with unique local solutions. Cover compact subintervals by finitely many chart pieces, glue by uniqueness and chart covariance using \cref{thm:pet-ch6-parallel-field-chart-local}, and exhaust the open interval by compact subintervals.
\end{proof}

\begin{theorem}[chart-local parallel fields]
  \label{thm:pet-ch6-parallel-field-chart-local}
  \lean{PetersenLib.parallelField_existence_uniqueness}
  \source{petersen:page-0252}
  \uses{def:pet-ch6-parallel-field}
  \leanok
  \dcref{ch6:6.1.3b}
  Let $c$ be a curve whose restriction to $(a,b)$ is differentiable in a chart
  around $\alpha$ and stays in that chart's source, with the associated
  Christoffel contraction continuous and bounded on $[a,b]$. Then for $t_0\in(a,b)$
  and $v\in T_{c(t_0)}M$ there is a vector field $V$ along $c$, parallel on
  $(a,b)$, with $V(t_0)=v$, and any parallel field $W$ along $c|_{(a,b)}$ with
  $W(t_0)=v$ agrees with $V$ on $(a,b)$.
\end{theorem}
\begin{proof}
  \uses{thm:pet-ch6-parallel-field-existence-uniqueness}
  \sketch
  Writing \(V=\sum_iV^iE_i\) gives \(\dot V=\sum_j(\dot V^j+\sum_i\alpha_i^jV^i)E_j\). The linear ODE existence and uniqueness theorem gives the claimed chart-local parallel field.
\end{proof}

\begin{definition}[Jacobi field]
  \label{def:pet-ch6-jacobi-field}
  \lean{PetersenLib.IsJacobiFieldAlong,PetersenLib.jacobiEquation,PetersenLib.isJacobiFieldAlong_add,PetersenLib.isJacobiFieldAlong_const_smul,PetersenLib.curveVelocity_isJacobiFieldAlong}
  \source{petersen:page-0253,petersen:page-0254}
  \uses{def:pet-ch6-vector-field-along-curve, lem:pet-ch6-third-partial-commutator}
  \notready
  Let $\bar c(s,t)$ be a \emph{geodesic variation} of a geodesic $c(t)=\bar
  c(0,t)$, i.e.\ $t\mapsto\bar c(s,t)$ is a geodesic for every $s$. By
  \cref{lem:pet-ch6-third-partial-commutator} applied with $u=t$ (using that
  $\frac{\partial^2\bar c}{\partial t^2}\equiv0$),
  \[
    0=\frac{\partial}{\partial s}\frac{\partial^2\bar c}{\partial t^2}
    =R\!\left(\frac{\partial c}{\partial s},\frac{\partial c}{\partial t}\right)
    \frac{\partial c}{\partial t}+\frac{\partial^2}{\partial t^2}\frac{\partial c}{\partial s},
  \]
  so the variational field $J(t)=\frac{\partial\bar c}{\partial s}(0,t)$ solves
  the linear second-order \emph{Jacobi equation} $\ddot J+R(J,\dot c)\dot c=0$.
  A \emph{Jacobi field} along $c$ is any solution of this equation; given
  $J(0),\dot J(0)$ there is a unique such solution (linearity of a second-order
  ODE). Conversely, every pair $(J(0),\dot J(0))$ with $J(0)=0$ arises from the
  geodesic variation $\bar c(s,t)=\exp_p(t(\dot c(0)+sJ(0)))$: since
  $\bar c(s,0)=p$ for all $s$, $J(0)=\frac{\partial\bar c}{\partial s}(0,0)=0$,
  and
  \[
    \dot J(0)=\frac{\partial^2\bar c}{\partial t\,\partial s}(0,0)
    =\frac{\partial^2\bar c}{\partial s\,\partial t}(0,0)
    =\frac{\partial}{\partial s}\big(\dot c(0)+sJ(0)\big)\Big|_{s=0}.
  \]

  \emph{Formalization note — what \texttt{\textbackslash leanok} covers here.} The marker
  certifies the \emph{definition}: \texttt{jacobiEquation} is $\ddot J+R(J,\dot c)\dot c$
  and \texttt{IsJacobiFieldAlong} is its vanishing. It does \emph{not} certify the
  existence-and-uniqueness theorem this node's exposition states in passing (a Jacobi field
  is determined by $(J(0),\dot J(0))$), which is not formalized. The converse realization of
  $J(0)=0$ data by the geodesic variation $\bar c(s,t)=\exp_p(t(\dot c(0)+sJ(0)))$ \emph{is}
  now formalized, in local form, by \cref{rem:pet-ch6-jacobi-dexp-relation}
  (\texttt{jacobiField\_dexp\_relation}).

  The Lean definition is chart-free and needs no chart hypothesis: $\ddot J$ is
  \texttt{derivAlongCurve} applied twice (that operator returns a field along $c$, and is
  chart-independent by \cref{def:pet-ch6-vector-field-along-curve}), and $R$ is the abstract
  \cref{def:pet-ch3-curvature-tensor}. Note the contrast with
  \cref{lem:pet-ch6-third-partial-commutator}, whose \emph{proof} needs a chart and
  \cref{lem:curvature-chart-bridge}.

  \emph{Naming.} The Lean name is \texttt{IsJacobiFieldAlong}, not
  \texttt{IsJacobiField}: the latter is already taken by \cref{def:pet-ch3-jacobi-field}
  for the \emph{distance-function} Jacobi field $L_{\partial_r}J=0$ of §3.2.4 — a different
  notion, with no geodesic and no second derivative. The suffix follows
  \cref{def:pet-ch6-parallel-field}'s \texttt{IsParallelAlong}.
\end{definition}

\begin{remark}[Jacobi fields and $D\exp_p$]
  \label{rem:pet-ch6-jacobi-dexp-relation}
  \lean{PetersenLib.jacobiField_dexp_relation,PetersenLib.expChart_radialVariation_deriv_eq,PetersenLib.expChart_radialVariation_hasDerivAt_on_ball,PetersenLib.expVariationField_eq_tangentCoordChange_dexp}
  \source{petersen:page-0254,petersen:page-0255}
  \uses{def:pet-ch6-jacobi-field}
  \notready
  For the geodesic variation $\bar c(s,t)=\exp_p(t(\dot c(0)+sJ(0)))$ of
  \cref{def:pet-ch6-jacobi-field},
  \[
    J(t)=\frac{\partial}{\partial s}\exp_p\big(t(\dot c(0)+sJ(0))\big)\Big|_{s=0}
    =D\exp_p\big(tJ(0)\big),
  \]
  where $tJ(0)\in T_{tv}T_pM$ for $v=\dot c(0)$. Consequently, $D\exp_p$ is
  nonsingular at $t_0v$ if and only if for every $J(t_0)\in
  T_{\exp_p(t_0v)}M$ there is a Jacobi field $J$ along $t\mapsto\exp_p(tv)$
  with $J(0)=0$ and value $J(t_0)$ at $t=t_0$.

  \emph{Formalization note.} The Lean statement is the local form Petersen's
  remark presumes throughout: it produces a normal radius $\rho$ and, for every
  $v$ with $\lvert v\rvert<\rho$ and every $w=J(0)$, a time window $(-\beta,\beta)$
  on which the variation field $\partial_s\bar c(\cdot,s)|_{s=0}$ of
  $\bar c(s,t)=\exp_p(t(v+sw))$ solves the Jacobi equation along
  $t\mapsto\exp_p(tv)$ and vanishes at $t=0$. This is the substantive content: the
  displayed identity $J=D\exp_p(t\,\cdot)$ merely names the variation field, since
  $\partial_s\exp_p(t(v+sw))|_{s=0}=D\exp_p|_{tv}(tw)$ by the chain rule (an
  equality the Lean does not separately restate, as it never mentions $D\exp_p$).
  The exponential rays $t\mapsto\exp_p(t(v+sw))$ are geodesics only while they stay
  in the normal ball, which is exactly what bounds the window. The final
  nonsingularity biconditional ($D\exp_p$ nonsingular at $t_0v$ iff every value
  $J(t_0)$ is realized) is the standard consequence for conjugate points and is not
  separately mechanized here.
\end{remark}

\begin{remark}[Jacobi fields and the Hessian of $r$]
  \label{rem:pet-ch6-jacobi-hessian-r}
  \lean{PetersenLib.jacobiField_hess_r}
  \source{petersen:page-0255}
  \uses{def:pet-ch6-jacobi-field}
  \leanok
  Let $c(t)$ be a unit speed geodesic with $c(0)=p$ and $J$ a Jacobi field
  along $c$ with $J(0)=0$; write $r(x)=|xp|$. As long as $tc'(0)$ lies in the
  interior of the domain on which $r$ is smooth (so that $\dot c(t)=\nabla
  r|_{c(t)}$), realizing $J$ via a geodesic variation $\bar c(s,t)$ as in
  \cref{def:pet-ch6-jacobi-field} gives
  \[
    \operatorname{Hess}r(J(t),J(t))=g(\nabla_{J(t)}\nabla r,J(t))
    =g\!\left(\frac{\partial^2\bar c}{\partial t\,\partial s},J\right)\Big|_{(0,t)}
    =g(\dot J(t),J(t)).
  \]
  The middle equality $g(\nabla_{J}\nabla r,J)=g(\dot J,J)$ is the crux: the
  Lean statement takes $J$ in the \S3.2.4 distance-function sense
  ($L_{\nabla r}J=0$), exactly the condition that turns
  $\nabla_{J}\nabla r$ into $\nabla_{\nabla r}J=\dot J$, and carries Petersen's
  radial hypothesis $\dot c(t)=\nabla r|_{c(t)}$ explicitly. The equation
  $\dot J=\nabla_{\dot c}J$ is the reusable along-curve $\leftrightarrow$ global
  Levi-Civita bridge , which
  identifies $\dot J(t)$ (the moving-foot covariant derivative
  \cref{def:pet-ch6-vector-field-along-curve}) with $\nabla_{\dot c(t)}J$ for a
  globally smooth field $J$; the first equality is the Ch.~3
  distance-function computation. The middle geometric interpretation via
  $\partial^2\bar c/\partial t\,\partial s$ is not separately restated, matching
  the honest scoping of \cref{rem:pet-ch6-jacobi-dexp-relation}.
\end{remark}

\begin{lemma}[the second variation integrand]
  \label{lem:pet-ch6-second-variation-integrand}
  \lean{PetersenLib.hasDerivAt_chartPairing_transversalAccel_of_geodesic_curvatureTensorAt,PetersenLib.hasDerivAt_chartPairing_transversalAccel_of_geodesic,PetersenLib.covariantDerivCoord_mixedPartial_fst_eq_sub_chartCurvature,PetersenLib.hasDerivAt_chartPairing_slice_ts,PetersenLib.surfaceCovariantDerivS_fderivSnd_eq_surfaceCovariantDerivT_fderivFst}
  \uses{def:pet-ch6-vector-field-along-curve}
  \leanok
  Let $\bar c(s,t)$ be a smooth variation whose base curve $t\mapsto\bar c(0,t)$ is
  a geodesic, and write $V=\partial c/\partial s$ and $T=\partial c/\partial t$.
  Then at $s=0$,
  \[
    \frac{\partial}{\partial t}g\Big(\frac{DV}{\partial s},T\Big)
      =g\Big(\frac{D}{\partial s}\frac{DT}{\partial s},T\Big)-g\big(R(V,T)V,T\big).
  \]
\end{lemma}
\begin{proof}
  \uses{lem:pet-ch6-third-partial-commutator}
  \sketch
  Metric compatibility, \(DT/\partial t=0\) on the geodesic base, equality \(DT/\partial s=DV/\partial t\), and the third-partial curvature commutator give the displayed integrand identity after pairing with \(T\).
\end{proof}


\begin{lemma}[the coordinate curvature in one chart, at every point of it]
  \label{lem:pet-ch6-offdiagonal-curvature-bridge}
  \lean{PetersenLib.curvatureTensorAt_chartBasis_of_mem,PetersenLib.curvatureTensor_coordinates_of_mem,PetersenLib.chartCurvatureContraction2_eq_neg_curvatureTensorAt_of_mem,PetersenLib.chartCurvature_eq_curvatureTensorAt_of_mem,PetersenLib.chartCurvature_coordChange}
  \source{petersen:page-0107}
  \uses{prop:pet-ch3-curvature-coordinate-formula}
  \leanok
  Fix a chart with domain $U$ and coordinate frame $\partial_1,\dots,\partial_n$.
  For \emph{every} $q\in U$ — not merely for one distinguished point of $U$ — the
  curvature tensor of the Levi-Civita connection satisfies
  \[
    R(\partial_i,\partial_j)\partial_k\big|_q=R^l_{ijk}(q)\,\partial_l\big|_q,
    \qquad
    R^l_{ijk}=\partial_i\Gamma^l_{jk}-\partial_j\Gamma^l_{ik}
      +\Gamma^s_{jk}\Gamma^l_{is}-\Gamma^s_{ik}\Gamma^l_{js},
  \]
  where $\Gamma^m_{ij}$ are the Christoffel symbols \emph{of that one chart},
  regarded as functions on $U$ and evaluated at $q$.

  Consequently the coordinate curvature of a chart computes the tensor $R$ at every
  point of the chart's domain, and is natural: if $q$ lies in the domains of two
  charts, their coordinate curvatures at $q$ agree after the change-of-frame
  isomorphism.
\end{lemma}
\begin{proof}
  \uses{prop:pet-ch3-curvature-coordinate-formula}
  \leanok
  \sketch
  Coordinate fields commute and satisfy \(\nabla_{\partial_i}\partial_j=\Gamma^m_{ij}\partial_m\). Leibniz expansion gives the four Christoffel terms at every chart point; tensoriality gives naturality under change of chart.
\end{proof}

\begin{theorem}[Synge's second variation formula]
  \label{thm:pet-ch6-synge-second-variation}
  \lean{PetersenLib.secondVariationEnergy,PetersenLib.secondVariationEnergy_deriv,PetersenLib.intervalIntegrable_syngeIntegrand_window,PetersenLib.syngeIntegrand_eq_chartCurvature,PetersenLib.hasDerivAt_deriv_windowEnergy,PetersenLib.hasDerivAt_deriv_windowEnergy_chart_curvatureTensorAt,PetersenLib.derivAlongCurve_variationField_eq_transfer,PetersenLib.energyFunctional_eventuallyEq_windowEnergy_chart,PetersenLib.secondVariationEnergy_chart_curvatureTensorAt,PetersenLib.secondVariationEnergy_chart,PetersenLib.hasDerivAt_chartPairing_transversalAccel_of_geodesic_curvatureTensorAt_of_mem,PetersenLib.iteratedDeriv_two_windowEnergy_chart,PetersenLib.hasDerivAt_deriv_windowEnergy_chart,PetersenLib.hasDerivAt_integral_chartPairing_ss,PetersenLib.hasDerivAt_chartPairing_slice_ss,PetersenLib.hasDerivAt_windowEnergy_chart_preByParts_shift}
  \source{petersen:page-0255}
  \uses{def:pet-ch6-vector-field-along-curve, lem:pet-ch6-offdiagonal-curvature-bridge}
  \leanok
  \dcref{ch6:6.1.4}
  If $\bar c:(-\varepsilon,\varepsilon)\times[a,b]\to M$ is a smooth variation
  of a geodesic $c(t)=\bar c(0,t)$, then
  \[
    \frac{d^2E(c_s)}{ds^2}\bigg|_{s=0}
    =\int_a^b\Big|\frac{\partial^2c}{\partial t\,\partial s}\Big|^2dt
    -\int_a^bg\!\left(R\!\left(\frac{\partial c}{\partial s},\frac{\partial c}{\partial t}\right)
    \frac{\partial c}{\partial t},\frac{\partial c}{\partial s}\right)dt
    +g\!\left(\frac{\partial^2c}{\partial s^2},\frac{\partial c}{\partial t}\right)\bigg|_a^b.
  \]
\end{theorem}
\begin{proof}
  \uses{lem:pet-ch6-third-partial-commutator,lem:pet-ch6-second-variation-integrand}
  \leanok
  \sketch
  Differentiate the energy twice on one chart slab, commute covariant derivatives with the curvature identity, and integrate the \(t\)-derivative to obtain the endpoint term. Cover the compact interval by slabs; chart-free integrands add and interior boundary terms telescope.
\end{proof}

\begin{remark}[special cases of the second variation formula]
  \label{rem:pet-ch6-second-variation-special-cases}
  \lean{PetersenLib.secondVariationEnergy_properVariation,PetersenLib.secondVariationEnergy_parallelField,PetersenLib.secondVariationEnergy_parallelField_geodesicTransversals,PetersenLib.curveAcceleration_const}
  \source{petersen:page-0256}
  \uses{thm:pet-ch6-synge-second-variation, def:pet-ch6-parallel-field}
  \leanok
  For a \emph{proper} variation (fixed endpoints) the boundary term in
  \cref{thm:pet-ch6-synge-second-variation} drops out, and with $V=\frac{\partial
  c}{\partial s}(0,\cdot)$,
  \[
    \frac{d^2E(c_s)}{ds^2}\bigg|_{s=0}=\int_a^b|\dot V|^2dt-\int_a^bg(R(V,\dot c)\dot c,V)dt.
  \]
  If instead $V$ is parallel, the first integral drops out:
  \[
    \frac{d^2E(c_s)}{ds^2}\bigg|_{s=0}=-\int_a^bg(R(V,\dot c)\dot c,V)dt
    +g\!\left(\frac{\partial^2c}{\partial s^2},\dot c\right)\bigg|_a^b,
  \]
  and if in addition $s\mapsto\bar c(s,t)$ is a geodesic for each $t$, the
  remaining boundary term also vanishes.
\end{remark}

\section{Nonpositive Sectional Curvature}

\begin{lemma}[nonsingular exponential maps are covering maps]
  \label{lem:pet-ch6-nonsingular-exp-covering}
  \lean{PetersenLib.exp_nonsingular_isCoveringMap,PetersenLib.exp_nonsingular_isCoveringMap_of_isProperMap,PetersenLib.exp_nonsingular_isCoveringMap_and_surjective}
  \source{petersen:page-0257}
  \dcref{ch6:6.2.1}
  If $\exp_p:T_pM\to M$ is nonsingular everywhere (no critical points), then it
  is a covering map.
\end{lemma}
\begin{proof}
  \sketch
  Nonsingular \(D\exp_p\) makes the pullback metric a local isometry on \(T_pM\). Radial lines are complete for this metric, so completeness of the domain gives the covering-map conclusion.
\end{proof}

\begin{theorem}[Hadamard--Cartan diffeomorphism theorem]
  \label{thm:pet-ch6-cartan-hadamard}
  \lean{PetersenLib.cartanHadamard,PetersenLib.cartanHadamard_homeomorphCore,PetersenLib.cartanHadamard_coveringCore,PetersenLib.jacobiField_eq_zero_on_Icc_of_endpoints_eq_zero_nonpositiveCurvature,PetersenLib.jacobiField_endpoint_ne_zero_nonpositiveCurvature}
  \notready
  \source{petersen:page-0257}
  \uses{lem:pet-ch6-nonsingular-exp-covering, def:pet-ch6-sec-bounds}
  \dcref{ch6:6.2.2}
  If $(M,g)$ is complete, connected, and has $\sec\le0$, then the universal
  cover of $M$ is diffeomorphic to $\mathbb{R}^n$.
\end{theorem}
\begin{proof}
  \uses{def:pet-ch6-jacobi-field, rem:pet-ch6-jacobi-dexp-relation, lem:pet-ch6-nonsingular-exp-covering}
  \sketch
  For \(J(0)=0,\dot J(0)=w\ne0\), \((|J|^2/2)''=|\dot J|^2-g(R(J,\dot c)\dot c,J)\ge|\dot J|^2\). Thus \(J(1)\ne0\), so \(\exp_p\) is nonsingular; \cref{lem:pet-ch6-nonsingular-exp-covering} gives the universal cover \(T_pM\cong\mathbb R^n\).
\end{proof}

\begin{remark}[failure for Ricci or scalar curvature bounds]
  \label{rem:pet-ch6-cartan-hadamard-fails-ricci}
  \lean{PetersenLib.cartanHadamard_fails_forRicci,PetersenLib.cartanHadamard_failure_of_notEuclidean}
  \notready
  \source{petersen:page-0258}
  \uses{thm:pet-ch6-cartan-hadamard}
  No analogue of \cref{thm:pet-ch6-cartan-hadamard} holds under $\operatorname{Ric}\le0$
  or $\operatorname{scal}\le0$ in place of $\sec\le0$: there exist Ricci-flat metrics on
  $\mathbb{R}^2\times S^{n-2}$ and scalar-flat metrics on $\mathbb{R}\times S^{n-1}$.
\end{remark}

\begin{definition}[strictly convex function]
  \label{def:pet-ch6-convex-function}
  \lean{PetersenLib.IsConvexOn,PetersenLib.IsStrictlyConvexOn}
  \source{petersen:page-0259}
  \leanok
  A function on a Riemannian manifold is \emph{(strictly) convex} if its
  restriction to every \emph{nonconstant} geodesic is (strictly) convex as a
  function of the affine parameter. The Lean statements are relative to an open
  set $U$ (``convex on $U$''), quantifying over the nonconstant geodesics that
  stay in $U$; this is what \Cref{thm:pet-ch6-convexity-radius-criterion}
  requires, and the global notion is the case $U=M$.

  Constant geodesics must be excluded: $f\circ c$ is then constant and cannot
  be strictly convex. The Lean predicate \texttt{IsNonconstantGeodesicOn}
  records this guard. The criterion used in
  \Cref{rem:pet-ch6-hessian-f0-positive-definite} is
  $(f_0\circ c)''=\operatorname{Hess}f_0(\dot c,\dot c)>0$.
\end{definition}

\begin{remark}[$\operatorname{Hess}f_0\succ0$ in nonpositive curvature]
  \label{rem:pet-ch6-hessian-f0-positive-definite}
  \lean{PetersenLib.hess_f0_posDef_nonpositiveCurvature}
  \source{petersen:page-0259}
  \uses{rem:pet-ch6-jacobi-hessian-r, def:pet-ch6-convex-function, thm:pet-ch6-cartan-hadamard}
  \notready
  Let $(M,g)$ be complete, simply connected, and of nonpositive curvature, so
  by \cref{thm:pet-ch6-cartan-hadamard} the distance $r(x)=|xp|$ is smooth on
  $M-\{p\}$ and $f_0(x)=\tfrac12r(x)^2$ is smooth everywhere, with
  $\operatorname{Hess}f_0=dr^2+r\operatorname{Hess}r$. If $J$ is a Jacobi field along a unit
  speed geodesic from $p$ with $J(0)=0$, then by
  \cref{rem:pet-ch6-jacobi-hessian-r} and the estimate in the proof of
  \cref{thm:pet-ch6-cartan-hadamard},
  \[
    \operatorname{Hess}r(J(b),J(b))=g(\dot J(b),J(b))\ge\int_0^b|\dot J|^2dt>0
  \]
  for $b>0$ and $J(b)\ne0$ (any $J(b)$ arises this way), so $\operatorname{Hess}r$, hence
  $\operatorname{Hess}f_0$, is positive definite. Consequently $f_0\circ c$ is (strictly)
  convex along any geodesic $c$, since
  $\frac{d^2}{dt^2}f_0\circ c=\operatorname{Hess}f_0(\dot c,\dot c)>0$; thus $f_0$ is a
  strictly convex function in the sense of
  \cref{def:pet-ch6-convex-function}.

  \emph{Formalization note.} The \texttt{\textbackslash leanok} content
  (\texttt{hess\_f0\_posDef\_nonpositiveCurvature}) is the genuinely-new analytic
  core: the \emph{strict positivity} $0<g(\dot J(b),J(b))$ of the displayed chain,
  for every along-curve Jacobi field ($\texttt{IsJacobiFieldAlong}$) with $J(0)=0$,
  $J(b)\ne0$, $b>0$, under $\sec\le0$ ($\texttt{HasSecBoundedAbove}\ g.\texttt{leviCivita}\ 0$).
  Via the already-\texttt{\textbackslash leanok} sibling
  \cref{rem:pet-ch6-jacobi-hessian-r} ($\operatorname{Hess}r(J,J)=g(\dot J,J)$) this
  reads $\operatorname{Hess}r(J(b),J(b))>0$. The proof runs entirely in the
  along-curve language: with $\varphi(t)=g(J,\dot J)$ one has
  $\varphi'=g(\dot J,\dot J)-g(R(J,\dot c)\dot c,J)\ge0$ (metric nonnegativity and
  $\sec\le0$, the denominator-cleared bound
  $\texttt{curvatureTensorFourAt\_le}$), so $\varphi$ is monotone; the fundamental
  theorem of calculus applied to $g(J,J)$ gives $\int_0^b2\varphi=|J(b)|^2>0$, which
  forces $\varphi(b)>0$. The middle inequality
  $g(\dot J(b),J(b))\ge\int_0^b|\dot J|^2$ — ``the estimate in the proof of
  \cref{thm:pet-ch6-cartan-hadamard}'' — is the separate lemma
  \texttt{jacobiField\_energy\_le\_boundary\_nonpositiveCurvature} (from
  \cref{rem:pet-ch6-ex-6-7-25}'s boundary identity plus $-g(R(J,\dot c)\dot c,J)\ge0$).
  The two remaining links to full positive-definiteness of $\operatorname{Hess}f_0$
  ---the surjectivity ``any $J(b)$ arises this way'' (\cref{thm:pet-ch6-cartan-hadamard}'s
  nonsingular $D\exp_p$) and the decomposition $\operatorname{Hess}f_0=dr^2+r\operatorname{Hess}r$---
  are the covering-theory pieces scoped out, exactly as in
  \cref{rem:pet-ch6-jacobi-hessian-r}.
\end{remark}

\begin{remark}[maxima of convex functions and unique minima]
  \label{rem:pet-ch6-max-of-convex-functions}
  \lean{PetersenLib.max_isConvex,PetersenLib.strictlyConvex_uniqueMinimum}
  \source{petersen:page-0259}
  \uses{def:pet-ch6-convex-function, thm:pet-ch5-hopf-rinow}
  \leanok
  The maximum of finitely many convex functions is convex (this reduces to
  the one-dimensional statement by restricting to geodesics). Any proper,
  nonnegative, strictly convex function $f$ on a complete manifold has a
  unique minimum: properness and boundedness below give a minimum, and if
  there were two minima, strict convexity restricted to a geodesic joining
  them would force smaller values on the interior of the segment than at
  either endpoint, a contradiction.

  \emph{Formalization notes.} (i) ``Maximum of finitely many'' is
  \texttt{Finset.sup'} over a \emph{nonempty} \texttt{Finset}: a maximum over
  an empty family has no value in $\mathbb{R}$. The $k$-point family of
  \cref{def:pet-ch6-linfty-center-of-mass} is nonempty, so this costs the
  consumer nothing. (ii) The Lean adds two hypotheses Petersen leaves
  implicit: $f$ is \emph{continuous} (properness alone does not give an
  attained infimum; Petersen's $f$ is a max of the smooth $\tfrac12r_{p_i}^2$,
  so this is free at the point of use, and deriving it from convexity would
  need convex~$\Rightarrow$~locally Lipschitz, which this project has not
  formalized), and $M$ is \emph{connected} (Petersen's standing assumption,
  and what \cref{thm:pet-ch5-hopf-rinow} needs to produce the joining
  geodesic). Properness is spelled as compact sublevel sets, which for
  nonnegative $f$ is equivalent to $f$ being a proper map. (iii)
  Nonnegativity is stated for fidelity but is not used --- compact sublevel
  sets already give existence.
\end{remark}

\begin{definition}[$L^\infty$ center of mass]
  \label{def:pet-ch6-linfty-center-of-mass}
  \lean{PetersenLib.centerOfMassLinfty}
  \source{petersen:page-0259,petersen:page-0260}
  \uses{rem:pet-ch6-hessian-f0-positive-definite, rem:pet-ch6-max-of-convex-functions}
  \leanok
  Let $(M,g)$ be complete, simply connected, of nonpositive curvature, and
  $p_1,\dots,p_k\in M$. By
  \cref{rem:pet-ch6-hessian-f0-positive-definite} and
  \cref{rem:pet-ch6-max-of-convex-functions}, the function
  $x\mapsto\max\{f_{0,p_1}(x),\dots,f_{0,p_k}(x)\}$ is proper,
  nonnegative, and strictly convex, hence has a unique minimum, denoted
  $\operatorname{cm}_\infty\{p_1,\dots,p_k\}$ and called the \emph{$L^\infty$ center of
  mass} of $\{p_1,\dots,p_k\}$: it is the center of the smallest closed ball
  containing $\{p_1,\dots,p_k\}$.
\end{definition}

\begin{theorem}[Cartan's center-of-mass theorem]
  \label{thm:pet-ch6-cartan-fixed-point}
  \lean{PetersenLib.cartan_finiteOrderIsometry_fixedPoint}
  \source{petersen:page-0260}
  \uses{def:pet-ch6-linfty-center-of-mass}
  \dcref{ch6:6.2.3}
  If $(M,g)$ is a complete simply connected Riemannian manifold of nonpositive
  curvature, then any isometry $F:M\to M$ of finite order has a fixed point.
\end{theorem}
\begin{proof}
  \sketch
  The finite orbit of \(p\) defines a proper strictly convex maximum of squared distances. \(F\) permutes the orbit, so its unique minimizer is fixed by \(F\).
\end{proof}

\begin{corollary}[torsion-freeness of the fundamental group]
  \label{cor:pet-ch6-torsion-free-fundamental-group}
  \lean{PetersenLib.fundamentalGroup_torsionFree_nonpositiveCurvature}
  \source{petersen:page-0261}
  \uses{thm:pet-ch6-cartan-fixed-point}
  \dcref{ch6:6.2.4}
  If $(M,g)$ is a complete Riemannian manifold of nonpositive curvature, then
  $\pi_1(M)$ is torsion free: every nontrivial element has infinite order.
\end{corollary}
\begin{proof}
  \uses{thm:pet-ch6-cartan-fixed-point}
  \sketch
  A finite-order deck transformation of the nonpositive-curvature universal cover has a fixed point by Cartan, contradicting the free deck action; hence \(\pi_1(M)\) is torsion free.
\end{proof}

\begin{lemma}[squared-distance Hessian lower bound]
  \label{lem:pet-ch6-hessian-lower-bound}
  \lean{PetersenLib.hess_f0_ge_metric_nonpositiveCurvature}
  \source{petersen:page-0261}
  \uses{rem:pet-ch6-hessian-f0-positive-definite}
  \dcref{ch6:6.2.5}
  If $(M,g)$ has nonpositive curvature, then any modified distance function
  $f_0=\tfrac12r^2$ satisfies $\operatorname{Hess}f_0\ge g$.
\end{lemma}
\begin{proof}
  \uses{rem:pet-ch6-jacobi-hessian-r}
  \sketch
  For a radial Jacobi field set \(\lambda=|J|^2/g(J,\dot J)\). Then \(\lambda(0)=0\) and the Jacobi equation, \(\sec\le0\), and Cauchy--Schwarz give \(\dot\lambda\le1\). Thus \(t\,\operatorname{Hess}r(J,J)\ge|J|^2\) and \(\operatorname{Hess}f_0\ge g\).
\end{proof}

\begin{remark}[triangle and quadrilateral comparison]
  \label{rem:pet-ch6-triangle-comparison}
  \lean{PetersenLib.triangleComparison_nonpositiveCurvature}
  \source{petersen:page-0262}
  \uses{lem:pet-ch6-hessian-lower-bound}
  Integrating $\frac{d}{dt}(f_{0,p}\circ c)\ge1$ along a unit speed geodesic
  $c$ (using \cref{lem:pet-ch6-hessian-lower-bound}) yields, for a triangle in
  $M$ with side lengths $a,b,c$ and angle $\alpha$ opposite $a$,
  \[
    a^2\ge b^2+c^2-2bc\cos\alpha.
  \]
  Consequently the angle sum in any geodesic triangle is $\le\pi$, and in any
  geodesic quadrilateral is $\le2\pi$. If $\sec<0$, these inequalities are
  strict unless the base point $p$ lies on the geodesic $c$; in particular the
  angle sum in any nondegenerate quadrilateral is $<2\pi$.
\end{remark}

\begin{definition}[axis, period, displacement function]
  \label{def:pet-ch6-axis-displacement}
  \lean{PetersenLib.IsAxis,PetersenLib.IsAxisPeriod,PetersenLib.axisPeriod, PetersenLib.displacementFunction}
  \source{petersen:page-0263}
  \leanok
  For an isometry $F:M\to M$, an \emph{axis} for $F$ is a geodesic
  $c:\mathbb{R}\to M$ such that $F\circ c$ is a reparametrization of $c$; since
  isometries map geodesics to geodesics, $F\circ c(t)=c(\pm t+a)$ for some
  $a\in\mathbb{R}$. If the sign is $-$, $c(a/2)$ is fixed by $F$. When
  $F\circ c(t)=c(t+a)$, $a$ is the \emph{period} of $F$ with respect to $c$
  (depending on the parametrization). The \emph{displacement function} of $F$
  is $x\mapsto\delta_F(x)=|xF(x)|$.

  Since Petersen's period is explicitly not unique, the relation
  ``$a$ is a period of $F$ with respect to $c$'' is primary and is the
  \texttt{Prop} \texttt{IsAxisPeriod}; \texttt{axisPeriod} is the
  derived $\mathbb{R}$-valued \emph{choice}, the infimum of the
  \emph{positive} periods. The positivity restriction is not cosmetic: over all
  periods the infimum would be junk on every periodic axis, since the period
  set is then $\{a_0+kP:k\in\mathbb{Z}\}$, unbounded below, on which
  $\operatorname{sInf}$ returns $0$ --- reporting ``$F$ fixes $c$ pointwise''
  for exactly the nontrivial translations of \Cref{thm:pet-ch6-preissmann}.
  Petersen's periods are positive in any case: \Cref{lem:pet-ch6-axis-existence}
  produces an axis of period $1$ from a unit-speed segment.
\end{definition}

\begin{lemma}[axis from positive minimal displacement]
  \label{lem:pet-ch6-axis-existence}
  \lean{PetersenLib.axis_of_positiveMinimalDisplacement}
  \source{petersen:page-0263}
  \uses{def:pet-ch6-axis-displacement}
  \dcref{ch6:6.2.7}
  If $F:M\to M$ is an isometry of a complete Riemannian manifold and
  $\delta_F$ has a positive minimum, then $F$ has an axis.
\end{lemma}
\begin{proof}
  \sketch
  At a minimizing point of the positive displacement, equality in the triangle inequality for \(c(t)\) and \(F(c(t))\) forces \(c\) followed by \(F\circ c\) to be geodesic. Iterating gives an axis of period \(1\).
\end{proof}

\begin{definition}[deck transformation]
  \label{def:pet-ch6-deck-transformation}
  \lean{PetersenLib.DeckTransformation,PetersenLib.DeckTransformation.isRiemannianIsometry}
  \source{petersen:page-0264}
  \leanok
  For the universal cover $\pi:\tilde M\to M$, a \emph{deck transformation} is
  a map $F:\tilde M\to\tilde M$ with $\pi\circ F=\pi$; it is determined by the
  value $F(p)\in\pi^{-1}(q)$ for a chosen $p\in\pi^{-1}(q)$, and the
  fundamental group $\pi_1(M,q)$ acts by deck transformations: a loop
  $[\alpha]\in\pi_1(M,q)$ and lift $\tilde\alpha$ with
  $\tilde\alpha(0)=p$ produce the deck transformation with
  $F(p)=\tilde\alpha(1)$. In the Riemannian setting, deck transformations are
  isometries of $\tilde M$ since $\pi$ is a local isometry.

  The Lean states this for an \emph{arbitrary} smooth covering $\pi:M\to M'$
  rather than specifically the universal cover --- the universal cover is
  available in neither Mathlib nor this project (see
  \Cref{thm:pet-ch6-cartan-hadamard}) --- and takes $F$ to be a diffeomorphism,
  which is automatic for deck transformations of a smooth covering. The
  operative closing claim, that deck transformations are isometries, is the
  proved \texttt{DeckTransformation.isRiemannianIsometry}, for the metric
  pulled back along $\pi$; it is the hinge \Cref{lem:pet-ch6-deck-transformation-dilation}
  needs. The remaining prose --- the $\pi_1(M,q)$-action and the ``determined by
  $F(p)$'' rigidity --- is \emph{not} formalized: both need covering-space
  lifting/monodromy machinery rather than Riemannian geometry.
\end{definition}

\begin{lemma}[minimal displacement of deck transformations]
  \label{lem:pet-ch6-deck-transformation-dilation}
  \lean{PetersenLib.deckTransformation_minimalDisplacement_closedGeodesic}
  \source{petersen:page-0264}
  \uses{def:pet-ch6-deck-transformation, def:pet-ch6-axis-displacement, lem:pet-ch6-axis-existence}
  \dcref{ch6:6.2.8}
  If $F:\tilde M\to\tilde M$ is a nontrivial deck transformation on the
  universal cover over a compact base $M$, then $\delta_F$ has a positive
  minimum; the axis through the minimizing point projects to a closed
  geodesic in $M$ of minimal length in its free homotopy class, and
  $\delta_F(x)\ge2\operatorname{inj}_{\pi(x)}(M)$.
\end{lemma}
\begin{proof}
  \sketch
  Projected paths from \(x\) to \(F(x)\) represent one free homotopy class, and nontriviality gives \(\delta_F(x)\ge2\operatorname{inj}_{\pi(x)}M\). Compactness of the base lets a minimizing sequence converge to a lifted path realizing the infimum; its projection is a shortest closed geodesic, and the axis lemma applies.
\end{proof}

\begin{theorem}[Preissmann's theorem]
  \label{thm:pet-ch6-preissmann}
  \lean{PetersenLib.preissmannTheorem,PetersenLib.preissmann_periods_injective_cyclic}
  \source{petersen:page-0263}
  \uses{lem:pet-ch6-deck-transformation-dilation, def:pet-ch6-axis-displacement, rem:pet-ch6-triangle-comparison}
  \dcref{ch6:6.2.6}
  If $(M,g)$ is a compact manifold of negative curvature, then every Abelian
  subgroup of $\pi_1(M)$ is cyclic. In particular, no compact product manifold
  $M\times N$ admits a metric of negative curvature.
\end{theorem}
\begin{proof}
  \uses{lem:pet-ch6-deck-transformation-dilation, rem:pet-ch6-triangle-comparison}
  \sketch
  By \cref{rem:pet-ch6-triangle-comparison}, negative curvature makes an axis unique through strict quadrilateral comparison. Commuting deck transformations preserve it, so periods embed an Abelian subgroup in \(\mathbb R\); the injectivity lower bound makes the image discrete and hence cyclic.
\end{proof}

\section{Positive Curvature}

\begin{definition}[sectional curvature bounds]
  \label{def:pet-ch6-sec-bounds}
  \lean{PetersenLib.HasSecBoundedBelow,PetersenLib.HasSecBoundedAbove, PetersenLib.HasSecIn,PetersenLib.HasSecPos,PetersenLib.HasSecNeg,PetersenLib.hasSecBoundedBelow_of_local_isometry_of_curvature,PetersenLib.hasSecBoundedAbove_of_local_isometry_of_curvature,PetersenLib.hasSecIn_of_local_isometry_of_curvature}
  \uses{def:pet-ch3-sectional-curvature}
  \leanok
  Let $(M,g)$ be a Riemannian manifold with Levi-Civita connection $D$ and let
  $k,K\in\mathbb{R}$. Say $(M,g)$ has $\sec\ge k$ if $\sec(v,w)\ge k$ for every
  $p\in M$ and every linearly independent pair $v,w\in T_pM$; $\sec\le K$ if
  $\sec(v,w)\le K$ for every such $p,v,w$; $k\le\sec\le K$ if both hold; and
  $\sec>0$ (resp.\ $\sec<0$) if $\sec(v,w)>0$ (resp.\ $<0$) for every such
  $p,v,w$. As $\sec$ depends only on the $2$-plane $\operatorname{span}\{v,w\}$,
  these are conditions on the $2$-planes of $TM$.

  This is the vocabulary of every comparison theorem of this section and of
  \emph{Basic Comparison Estimates} below, and it is stated here because
  Petersen uses the notation throughout without ever displaying it as a
  definition.

  Two points are forced rather than chosen. The linear independence hypothesis
  is essential: $\sec(v,w)$ is a quotient by $g(v\wedge w,v\wedge w)$, which
  vanishes exactly on dependent pairs, so dropping it would make $\sec\ge k$
  unsatisfiable for every $k>0$ and every downstream theorem vacuous. And
  $\sec>0$ is \emph{not} $\exists k>0,\ \sec\ge k$: on a noncompact manifold
  positive curvature may decay to $0$ at infinity (a paraboloid) with no uniform
  positive lower bound, so \Cref{thm:pet-ch6-synge-theorem} needs its own
  predicate.
\end{definition}

\begin{definition}[Ricci curvature lower bound]
  \label{def:pet-ch6-ricci-lower-bound}
  \lean{PetersenLib.HasRicciBoundedBelowAt,PetersenLib.HasRicciBoundedBelow}
  \uses{def:pet-ch3-ricci-curvature}
  \leanok
  Let $n=\dim M$.  Say that $(M,g)$ has
  $\operatorname{Ric}\ge(n-1)k\,g$ if, for every $p\in M$ and $v\in T_pM$,
  \[
    (n-1)k\,g(v,v)\le \operatorname{Ric}(v,v).
  \]
  In particular, for a unit vector $v$ this reads
  $\operatorname{Ric}(v,v)\ge(n-1)k$.
\end{definition}

\begin{lemma}[Bonnet--Synge index inequality, scalar core]
  \label{lem:pet-ch6-bonnet-synge-index-core}
  \lean{PetersenLib.bonnetSynge_index_core,PetersenLib.integral_sin_sq_window, PetersenLib.integral_cos_sq_window,PetersenLib.sq_pi_div_lt_of_pi_div_sqrt_lt}
  \source{petersen:page-0266,petersen:page-0267}
  \uses{def:pet-ch6-sec-bounds}
  \leanok
  Fix $l>0$, $k>0$ with $l>\pi/\sqrt k$, and let $\kappa:[0,l]\to\mathbb R$ satisfy
  $\kappa\ge k$ on $[0,l]$, with $\sin^2(\pi\cdot/l)\,\kappa$ integrable. Then
  \[
    \Big(\frac\pi l\Big)^2\int_0^l\cos^2\Big(\frac\pi lt\Big)dt
      -\int_0^l\sin^2\Big(\frac\pi lt\Big)\kappa(t)\,dt<0 .
  \]
  This is the pure-analysis heart of Bonnet--Synge (and, summed over an orthonormal frame,
  of Myers, \cref{thm:pet-ch6-myers-ricci-diameter}); the ingredients are the window
  integrals $\int_0^l\sin^2(\pi t/l)\,dt=\int_0^l\cos^2(\pi t/l)\,dt=l/2$ and the equivalence
  of $l>\pi/\sqrt k$ with $(\pi/l)^2<k$.
\end{lemma}
\begin{proof}
  \leanok
  \sketch
  Both trigonometric integrals equal \(l/2\). Since \(\kappa\ge k\) and \((\pi/l)^2<k\), the expression is at most \((l/2)((\pi/l)^2-k)<0\).
\end{proof}

\begin{lemma}[Bonnet--Synge index inequality]
  \label{lem:pet-ch6-bonnet-synge-diameter}
  \lean{PetersenLib.bonnetSynge_longGeodesicsNotMinimizing, PetersenLib.bonnetSynge_longGeodesicsNotMinimizing_of_secLowerBound, PetersenLib.bonnetSynge_distance_lt_of_secLowerBound, PetersenLib.bonnetSynge_variation_integrand}
  \source{petersen:page-0266}
  \uses{rem:pet-ch6-second-variation-special-cases, def:pet-ch6-sec-bounds, %
    lem:pet-ch6-bonnet-synge-index-core}
  \notready
  \dcref{ch6:6.3.1}
  If $(M,g)$ satisfies $\sec\ge k>0$, then geodesics of length $>\pi/\sqrt k$
  cannot be locally minimizing.
\end{lemma}
\begin{proof}
  \uses{rem:pet-ch6-second-variation-special-cases,lem:pet-ch6-bonnet-synge-index-core}
  \sketch
  Use \(V(t)=\sin(\pi t/l)E(t)\) in the proper second-variation formula. The preceding scalar inequality makes the second variation negative when \(l>\pi/\sqrt{k}\), so the geodesic is not minimizing.
\end{proof}

\begin{corollary}[Hopf--Rinow and Myers compactness criterion]
  \label{cor:pet-ch6-hopf-rinow-myers-diameter}
  \lean{PetersenLib.hopfRinowMyers_finiteFundamentalGroup_of_explicitCover}
  \source{petersen:page-0267}
  \uses{lem:pet-ch6-bonnet-synge-diameter, def:pet-ch6-sec-bounds}
  \notready
  \dcref{ch6:6.3.2}
  \notready
  If $(M,g)$ is complete and satisfies $\sec\ge k>0$, then $M$ is compact and
  $\operatorname{diam}(M,g)\le\pi/\sqrt k=\operatorname{diam}S^n_k$; in particular $\pi_1(M)$ is
  finite.
\end{corollary}
\begin{proof}
  \uses{lem:pet-ch6-bonnet-synge-diameter}
  \sketch
  By \cref{lem:pet-ch6-bonnet-synge-diameter}, no minimizing geodesic exceeds \(\pi/\sqrt{k}\). Hopf--Rinow gives the diameter bound and compactness; the universal cover is compact, so its deck group is finite.
\end{proof}

\begin{corollary}[Bonnet--Myers diameter and compactness]
  \label{cor:pet-ch6-hopf-rinow-myers-geometric}
  \lean{PetersenLib.hopfRinowMyers_diameterBound_of_secLowerBound}
  \source{petersen:page-0267}
  \uses{def:pet-ch6-sec-bounds, thm:pet-ch6-myers-ricci-geometric}
  \leanok
  Let $(M^n,g)$ be complete and connected, with $n\ge2$. If
  $\sec\ge k>0$, then $M$ is compact and
  $\operatorname{diam}(M,g)\le\pi/\sqrt k$.
\end{corollary}
\begin{proof}
  \leanok
  \sketch
  Tracing \(\sec\ge k\) over an orthonormal basis gives \(\operatorname{Ric}\ge(n-1)k\); apply \cref{thm:pet-ch6-myers-ricci-geometric}.
\end{proof}

\begin{lemma}[Myers' second-variation core]
  \label{lem:pet-ch6-myers-index}
  \lean{PetersenLib.myersRicci_secondVariation_neg,PetersenLib.bonnetSynge_secondVariation_eq, PetersenLib.myers_index_core, PetersenLib.ricciCurvature_eq_sum_sectional_of_velocitySeededFrame}
  \source{petersen:page-0267,petersen:page-0268}
  \uses{lem:pet-ch6-bonnet-synge-index-core, rem:pet-ch6-second-variation-special-cases, def:pet-ch6-sec-bounds}
  \leanok
  Let $c:[0,l]\to M$ be a unit-speed geodesic with $l>\pi/\sqrt k$, $k>0$, and let
  $E_1=\dot c,E_2,\dots,E_n$ be a $g$-orthonormal parallel frame along $c$ with
  $\operatorname{Ric}(\dot c,\dot c)=\sum_{i=2}^n\sec(E_i,\dot c)\ge(n-1)k$. Then, taking the
  $m=n-1$ perpendicular fields $V_i(t)=\sin(\pi t/l)E_i(t)$, at least one has a strictly
  negative second variation $\tfrac{d^2}{ds^2}E(f_{i,s})|_0<0$; hence $c$ is not locally
  minimizing.

  \emph{Formalization note.} This is the analytic heart of \cref{thm:pet-ch6-myers-ricci-diameter}.
  Each perpendicular field gives second variation
  $\tfrac{d^2}{ds^2}E(f_{i,s})|_0=(\pi/l)^2(l/2)-\int_0^l\sin^2(\pi t/l)\sec(E_i,\dot c)\,dt$
  (\texttt{bonnetSynge\_secondVariation\_eq}, the \cref{lem:pet-ch6-bonnet-synge-diameter}
  computation stopped one step early to keep the value); summing over the frame and using the
  frame-sum Ricci bound $(n-1)k\le\sum_i\sec(E_i,\dot c)$ makes the total strictly negative
  (\texttt{myers\_index\_core}), so some summand is $<0$. As in
  \cref{lem:pet-ch6-bonnet-synge-diameter} the variations, their slab-smoothness, the frame's
  parallelism/orthonormality, and the integrability are hypotheses. The Ricci bound is carried in
  its frame-sum form; it equals $\operatorname{Ric}(\dot c,\dot c)\ge(n-1)k$ by the trace identity
  $\operatorname{Ric}(\dot c,\dot c)=\sum_i\sec(E_i,\dot c)$
  (\texttt{ricciCurvature\_eq\_sum\_sectionalCurvature}, §3.1.4) for the orthonormal frame.
\end{lemma}
\begin{proof}
  \uses{lem:pet-ch6-bonnet-synge-index-core,rem:pet-ch6-second-variation-special-cases}
  \leanok
  \sketch
  Sum the second variations of \(V_i=\sin(\pi t/l)E_i\) over \(i=2,\ldots,n\) using \cref{rem:pet-ch6-second-variation-special-cases}. The Ricci trace and \((\pi/l)^2<k\) make the sum negative, so one field gives a negative variation.
\end{proof}

\begin{theorem}[Myers diameter and compactness]
  \label{thm:pet-ch6-myers-ricci-geometric}
  \lean{PetersenLib.myersRicciDiameterBound_of_ricciLowerBound}
  \source{petersen:page-0267}
  \uses{lem:pet-ch6-myers-index, def:pet-ch6-ricci-lower-bound}
  \leanok
  Let $(M^n,g)$ be complete and connected, with $n\ge2$. If
  $\operatorname{Ric}\ge(n-1)k$ for some $k>0$, then $M$ is compact and
  $\operatorname{diam}(M,g)\le\pi/\sqrt k$.
\end{theorem}
\begin{proof}
  \leanok
  \sketch
  The estimate in \cref{lem:pet-ch6-myers-index} excludes minimizing geodesics longer than \(\pi/\sqrt{k}\); Hopf--Rinow yields the diameter bound and compactness.
\end{proof}

\begin{theorem}[Myers on an explicit simply connected cover]
  \label{thm:pet-ch6-myers-explicit-cover}
  \lean{PetersenLib.myersRicci_of_explicit_simplyConnectedCover}
  \uses{thm:pet-ch6-myers-ricci-geometric}
  \leanok
  Let $\widetilde M^n$ be a complete, connected, simply connected Riemannian
  manifold with $n\ge2$ and $\operatorname{Ric}_{\widetilde M}\ge(n-1)k$ for
  some $k>0$. If $p:\widetilde M\to M$ is a surjective covering map onto a
  $T_1$ space, then
  $\operatorname{diam}(\widetilde M)\le\pi/\sqrt k$, both $\widetilde M$ and
  $M$ are compact, and $\pi_1(M,y)$ is finite for every $y\in M$.
\end{theorem}
\begin{proof}
  \leanok
  \sketch
  Apply Myers to the simply connected cover, obtaining compactness. Its image is compact and each fiber is closed discrete, hence finite; monodromy injects \(\pi_1(M)\) into a finite fiber.
\end{proof}

\begin{theorem}[Myers' fundamental-group theorem]
  \label{thm:pet-ch6-myers-ricci-diameter}
  \source{petersen:page-0267}
  \uses{lem:pet-ch6-myers-index, cor:pet-ch6-hopf-rinow-myers-diameter, def:pet-ch6-ricci-lower-bound}
  \notready
  \dcref{ch6:6.3.3}
  \notready
  If $(M,g)$ is a complete Riemannian manifold with $\operatorname{Ric}\ge(n-1)k>0$,
  then $\operatorname{diam}(M,g)\le\pi/\sqrt k$, and $(M,g)$ has finite fundamental
  group.
\end{theorem}
\begin{proof}
  \uses{rem:pet-ch6-second-variation-special-cases}
  \sketch
  The summed sine-field variation contradicts any minimizing geodesic longer than \(\pi/\sqrt{k}\). The diameter bound follows, and the compact universal cover has a finite deck group.
\end{proof}

\begin{example}[the punctured sphere]
  \label{ex:pet-ch6-punctured-sphere-infinite-pi1}
  \lean{PetersenLib.puncturedSphereExample}
  \notready
  \source{petersen:page-0268}
  \dcref{ch6:6.3.4}
  The incomplete manifold $S^2-\{\pm p\}$ has constant curvature $1$ and
  infinite fundamental group; its universal cover also has diameter $\pi$.
  This shows completeness cannot be dropped in
  \cref{cor:pet-ch6-hopf-rinow-myers-diameter} or
  \cref{thm:pet-ch6-myers-ricci-diameter}.
\end{example}

\begin{example}[a doubly warped product with positive Ricci curvature]
  \label{ex:pet-ch6-doubly-warped-ricci-positive}
  \lean{PetersenLib.dwpRicciPositiveOnS1TimesR3}
  \notready
  \source{petersen:page-0268}
  \uses{thm:pet-ch6-myers-ricci-diameter}
  \dcref{ch6:6.3.5}
  $S^1\times\mathbb{R}^3$ admits a complete doubly warped product metric
  $dr^2+\rho^2(r)d\theta^2+\phi^2(r)ds_r^2$ with $\operatorname{Ric}>0$ everywhere,
  showing incompleteness is not needed to violate the topological conclusions
  of \cref{thm:pet-ch6-myers-ricci-diameter} on noncompact manifolds — Myers'
  theorem itself forces compactness once $\operatorname{Ric}$ is bounded below by a
  positive constant, but here $\operatorname{Ric}$ is merely positive (not bounded
  away from $0$). With $\rho(t)=t^{-1/4}$, $\phi(t)=t^{3/4}$ for $t>1$ the
  Ricci curvature is positive; extending to $[0,\infty)$ so the metric is
  $C^\infty$ at $0$, one may arrange $\rho,\phi$ to be $C^1$ and piecewise
  smooth at $t=1$ with $-1<\dot\rho\le0$, $0<\dot\phi\le1$, and, on $[0,1]$,
  $\dot\rho\le0$; smoothing $\rho$ at $t=1$ while preserving positivity of
  $\operatorname{Ric}$ then yields the example.
\end{example}

\begin{remark}[closed parallel fields around closed geodesics]
  \label{rem:pet-ch6-closed-parallel-field-construction}
  \lean{PetersenLib.closedParallelFieldAroundClosedGeodesic,PetersenLib.closedParallelField_holonomyFixedVector}
  \notready
  \source{petersen:page-0268,petersen:page-0269}
  \uses{def:pet-ch6-parallel-field}
  Let $c:[0,l]\to M$ be a closed unit speed geodesic, $p=c(0)=c(l)$, and
  $P:T_pM\to T_pM$ parallel translation once around $c$, a linear isometry
  fixing $\dot c(0)$ (since $P(\dot c(0))=\dot c(l)=\dot c(0)$), hence
  preserving $\dot c(0)^\perp\subset T_pM$. A linear isometry
  $L:\mathbb{R}^k\to\mathbb{R}^k$ with $\det L=(-1)^{k+1}$ has $1$ as an eigenvalue.
  Consequently: (1) if $M$ is orientable and even-dimensional, $P$ preserves
  orientation ($\det P=1$) and $\dot c(0)^\perp$ is odd-dimensional, so
  $P|_{\dot c(0)^\perp}$ has $\det=(-1)^{(\dim M-1)+1}=1$... more precisely,
  since $\dim(\dot c(0)^\perp)=\dim M-1$ is odd, $P|_{\dot c(0)^\perp}$ fixes
  a vector, giving a closed parallel field around $c$; (2) if $M$ is
  nonorientable, odd-dimensional, and $c$ is orientation-reversing, then
  $\dot c(0)^\perp$ is even-dimensional and $P|_{\dot c(0)^\perp}$ has
  $\det=-1$, again fixing a vector and yielding a closed parallel field.
\end{remark}

\begin{theorem}[Synge's orientability theorem]
  \label{thm:pet-ch6-synge-theorem}
  \lean{PetersenLib.syngeTheorem}
  \notready
  \source{petersen:page-0269}
  \uses{lem:pet-ch6-deck-transformation-dilation, rem:pet-ch6-closed-parallel-field-construction, rem:pet-ch6-second-variation-special-cases, def:pet-ch6-sec-bounds}
  \dcref{ch6:6.3.6}
  Let $M$ be a compact manifold with $\sec>0$. (1) If $M$ is even-dimensional
  and orientable, $M$ is simply connected. (2) If $M$ is odd-dimensional,
  $M$ is orientable.
\end{theorem}
\begin{proof}
  \uses{lem:pet-ch6-deck-transformation-dilation, rem:pet-ch6-closed-parallel-field-construction, rem:pet-ch6-second-variation-special-cases}
  \sketch
  For the even-dimensional assertion, suppose that an orientable $M$ has a
  nontrivial deck transformation $F$. For the odd-dimensional assertion,
  suppose that $M$ is nonorientable and choose $F$ representing an
  orientation-reversing loop. In either case an axis of $F$ projects to a
  shortest closed geodesic $c$ in its free homotopy class. In the first case
  the normal holonomy acts on an odd-dimensional space and preserves
  orientation; in the second it acts on an even-dimensional space and reverses
  orientation. Thus \cref{rem:pet-ch6-closed-parallel-field-construction}
  supplies a closed unit parallel field $E\perp\dot c$ in precisely the two
  required parity cases. The closed variation generated by $E$ has
  \[
    \frac{d^2}{ds^2}E(c_s)\bigg|_{s=0}
      =-\int_c\sec(E,\dot c)\,dt<0
  \]
  by \cref{rem:pet-ch6-second-variation-special-cases}, contradicting the
  minimality of $c$. Hence the even orientable case is simply connected and an
  odd-dimensional positively curved manifold cannot be nonorientable.
\end{proof}

\begin{remark}[the Hopf problem]
  \label{rem:pet-ch6-hopf-problem}
  \lean{PetersenLib.hopfProblemRemark}
  \source{petersen:page-0269}
  \uses{thm:pet-ch6-synge-theorem}
  \Cref{thm:pet-ch6-synge-theorem} shows $\mathbb{RP}^2\times\mathbb{RP}^2$, which has
  positive Ricci curvature, admits no metric of positive sectional curvature
  (it is even-dimensional, orientable, but not simply connected). Whether
  $S^2\times S^2$ admits positive sectional curvature is open (the Hopf
  problem); \cref{thm:pet-ch6-preissmann} ruled out negative curvature on any
  product manifold via fundamental group considerations, but such
  considerations alone do not resolve the positive-curvature case.
\end{remark}

\section{Basic Comparison Estimates}

\begin{proposition}[Riccati comparison principle]
  \label{prop:pet-ch6-riccati-comparison-principle}
  \lean{PetersenLib.riccatiComparisonPrinciple}
  \source{petersen:page-0270}
  \leanok
  \dcref{ch6:6.4.1a}
  Let $\rho_1,\rho_2:(0,b)\to\mathbb{R}$ be smooth with
  \[
    \dot\rho_1+\rho_1^2\le\dot\rho_2+\rho_2^2 ,
  \]
  and let $F$ be an antiderivative of $\rho_1+\rho_2$ on $(0,b)$. Then
  $(\rho_2-\rho_1)e^{F}$ is nondecreasing on $(0,b)$; equivalently, for all
  $0<s\le t<b$,
  \[
    \rho_2(t)-\rho_1(t)\ \ge\ \big(\rho_2(s)-\rho_1(s)\big)
      \exp\Big({-\int_s^t(\rho_1+\rho_2)}\Big).
  \]
\end{proposition}
\begin{proof}
  \leanok
  \sketch
  For \(G=(\rho_2-\rho_1)e^F\), \(G'=(\dot\rho_2-\dot\rho_1+\rho_2^2-\rho_1^2)e^F\ge0\). Integration gives the exponential comparison inequality.
\end{proof}

\begin{remark}[correction to the Riccati conclusion]
  \label{rem:pet-ch6-riccati-principle-correction}
  \lean{PetersenLib.riccatiComparisonPrinciple_statement_counterexample}
  \source{petersen:page-0270}
  \uses{prop:pet-ch6-riccati-comparison-principle}
  \leanok
  \dcref{ch6:6.4.1b}
  The stronger endpoint assertion
  \[
    \rho_2-\rho_1\ \ge\ \limsup_{t\to0}\big(\rho_2(t)-\rho_1(t)\big),
  \]
  which is false as stated. On $(0,2)$ put
  \[
    \rho_1(t)=\frac{1-e^{-t}}{2},\qquad \rho_2(t)=\frac{1+e^{-t}}{2},
  \]
  so $\rho_1+\rho_2=1$ and $\rho_2-\rho_1=e^{-t}$. Then
  $\rho_2^2-\rho_1^2=(\rho_2-\rho_1)(\rho_2+\rho_1)=e^{-t}$ while
  $\dot\rho_2-\dot\rho_1=-e^{-t}$, so $\dot\rho_1+\rho_1^2=\dot\rho_2+\rho_2^2$
  and the hypothesis holds with equality. Yet
  $\limsup_{t\to0}(\rho_2-\rho_1)=1$ while $\rho_2(1)-\rho_1(1)=e^{-1}<1$; since
  $\rho_2-\rho_1=e^{-t}$ is strictly decreasing, the asserted inequality in fact
  fails at \emph{every} $t\in(0,2)$.

  Monotonicity establishes only that
  $(\rho_2-\rho_1)e^{F}$ is nondecreasing. Deducing the asserted inequality from
  that would require $e^{F(s)-F(t)}\ge1$ in the limit $s\to0^+$, which is not
  automatic: here $F(t)=t$, so $e^{F(s)-F(t)}\to e^{-t}<1$. The monotonicity
  itself is correct, and it is all that
  \cref{cor:pet-ch6-riccati-comparison-estimate} consumes — through
  \cref{lem:pet-ch6-riccati-comparison-le} — so no downstream result is affected.
\end{remark}

\begin{lemma}[Riccati comparison — corrected conclusion]
  \label{lem:pet-ch6-riccati-comparison-le}
  \lean{PetersenLib.riccatiComparisonPrinciple_le_of_liminf}
  \uses{prop:pet-ch6-riccati-comparison-principle}
  \leanok
  In the setting of \cref{prop:pet-ch6-riccati-comparison-principle}, suppose in
  addition that
  \[
    \liminf_{t\to0^+}\big(\rho_2(t)-\rho_1(t)\big)e^{F(t)}\ \ge\ 0 .
  \]
  Then $\rho_1\le\rho_2$ on all of $(0,b)$.
\end{lemma}
\begin{proof}
  \leanok
  \uses{prop:pet-ch6-riccati-comparison-principle}
  \sketch
  The preceding monotonicity and the nonnegative liminf at \(0^+\) imply \(G\ge0\), hence \(\rho_1\le\rho_2\).
\end{proof}

\begin{remark}[a sufficient initial condition for Riccati comparison]
  \label{rem:pet-ch6-riccati-liminf-hypothesis}
  \lean{PetersenLib.riccatiComparisonPrinciple_le_of_bounded_of_exp_tendsto_zero}
  \source{petersen:page-0270}
  \uses{lem:pet-ch6-riccati-comparison-le}
  \leanok
  If $\rho_2-\rho_1$ is bounded near $0^+$ and $e^{F}\to0$ as $t\to0^+$, then the
  side condition of \cref{lem:pet-ch6-riccati-comparison-le} holds, and hence
  $\rho_1\le\rho_2$ on $(0,b)$. The second hypothesis amounts to
  $\int_{0^+}(\rho_1+\rho_2)=-\infty$. This is the situation of
  \cref{cor:pet-ch6-riccati-comparison-estimate}, where
  $\rho_1(t),\rho_2(t)=\tfrac1t+O(t)$: then $\rho_1+\rho_2=\tfrac2t+O(t)$ gives
  $F(t)=2\log t+O(t^2)$ and $e^{F(t)}=t^2e^{O(t^2)}$, while $\rho_2-\rho_1=O(t)$,
  so $G=(\rho_2-\rho_1)e^{F}=O(t^3)\to0$ and the hypothesis holds. In that regime
  $\limsup_{t\to0}(\rho_2-\rho_1)=0$, so the endpoint assertion reduces to
  $\rho_1\le\rho_2$.

  Both halves of the hypothesis are needed on each $\rho_i$ separately, not just
  on the sum: for $\rho_{1,2}(t)=\tfrac1t\mp\tfrac1{2t^2}$ one still has
  $\rho_1+\rho_2=\tfrac2t$ and $\dot\rho_1+\rho_1^2=\dot\rho_2+\rho_2^2$, but
  $\rho_2-\rho_1=t^{-2}$ is unbounded near $0$ and
  $\limsup_{t\to0}(\rho_2-\rho_1)=+\infty$.
\end{remark}
\begin{proof}
  \leanok
  \uses{lem:pet-ch6-riccati-comparison-le}
  \sketch
  Bounded \(\rho_2-\rho_1\) and \(e^F\to0\) force \(G=(\rho_2-\rho_1)e^F\) to have nonnegative liminf, so \cref{rem:pet-ch6-riccati-liminf-hypothesis} applies.
\end{proof}

\begin{corollary}[Riccati comparison estimate]
  \label{cor:pet-ch6-riccati-comparison-estimate}
  \lean{PetersenLib.riccatiComparisonEstimate}
  \source{petersen:page-0270,petersen:page-0271}
  \uses{prop:pet-ch6-riccati-comparison-principle, lem:pet-ch6-riccati-comparison-le, rem:pet-ch6-riccati-liminf-hypothesis, rem:pet-ch6-riccati-estimate-strictness}
  \leanok
  \dcref{ch6:6.4.2a}
  Let $\rho:(0,b)\to\mathbb{R}$ be smooth with $\rho(t)=\tfrac1t+O(t)$ and $k\in\mathbb{R}$.
  (1) If $\dot\rho+\rho^2\le-k$, then $\rho(t)\le\operatorname{sn}_k'(t)/\operatorname{sn}_k(t)$, and
  moreover $b\le\pi/\sqrt k$ when $k>0$. (2) If $-k\le\dot\rho+\rho^2$, then
  $\operatorname{sn}_k'(t)/\operatorname{sn}_k(t)\le\rho(t)$ for all $t<b$ when $k\le0$, and for
  $t<\min\{b,\pi/\sqrt k\}$ when $k>0$.

  The endpoint is non-strict; see
  \cref{rem:pet-ch6-riccati-estimate-strictness}.
\end{corollary}

\begin{remark}[the non-strict Riccati endpoint bound]
  \label{rem:pet-ch6-riccati-estimate-strictness}
  \lean{PetersenLib.riccatiComparisonEstimate_lt_isFalse}
  \source{petersen:page-0271}
  \leanok
  \dcref{ch6:6.4.2b}
  The bound $b\le\pi/\sqrt k$ in \cref{cor:pet-ch6-riccati-comparison-estimate}(1)
  cannot be strengthened to $b<\pi/\sqrt k$. Take $k=1$ and
  $\rho=\operatorname{sn}_1'/\operatorname{sn}_1=\cot$ on $b=\pi=\pi/\sqrt k$. Then $\rho$ is
  smooth on $(0,\pi)$, satisfies $\rho(t)=\tfrac1t+O(t)$, and solves
  $\dot\rho+\rho^2=-1$, so every hypothesis of (1) holds — with equality — yet
  $b=\pi/\sqrt k$. The printed argument rules out $b>\pi/\sqrt k$ (there $\rho$
  would have to blow up before $\pi/\sqrt k$, contradicting smoothness on
  $(0,b)$), which gives exactly $b\le\pi/\sqrt k$; the boundary case is realised,
  as the model $\rho_k$ itself shows.

  The model case records the sharp endpoint and prevents replacing $\le$ by a
  strict inequality in later uses.
\end{remark}
\begin{proof}
  \uses{lem:pet-ch6-riccati-comparison-le,rem:pet-ch6-riccati-liminf-hypothesis}
  \sketch
  The model \(\rho_k=\operatorname{sn}_k'/\operatorname{sn}_k\) has first pole at \(\pi/\sqrt{k}\) for \(k>0\). Comparison rules out larger \(b\), while \(\rho=\cot t\) on \(b=\pi\) realizes equality.
\end{proof}

\begin{theorem}[Rauch comparison theorem]
  \label{thm:pet-ch6-rauch-comparison}
  \lean{PetersenLib.rauchComparisonHessian,PetersenLib.rauchComparisonHessian_radialField}
  \source{petersen:page-0271}
  \uses{cor:pet-ch6-riccati-comparison-estimate, rem:pet-ch6-jacobi-hessian-r, def:pet-ch6-sec-bounds}
  \dcref{ch6:6.4.3}
  Assume $(M,g)$ satisfies $k\le\sec\le K$, and write $g=dr^2+g_r$ for the
  metric in exponential polar coordinates around $p$ on the star-shaped
  region $T_pM-\{0\}$. Then
  \[
    \frac{\operatorname{sn}_K'(r)}{\operatorname{sn}_K(r)}g_r\le\operatorname{Hess}r\le\frac{\operatorname{sn}_k'(r)}{\operatorname{sn}_k(r)}g_r,
  \]
  and consequently the modified distance functions $f_k,f_K$ satisfy
  $\operatorname{Hess}f_k\le(1-kf_k)g$ and $\operatorname{Hess}f_K\ge(1-Kf_K)g$.

  The ordering of the comparison functions is essential: the upper curvature
  bound $K$ gives the lower Hessian bound and the lower curvature bound $k$ gives
  the upper one, as follows from
  $\dot\rho\ge-K-\rho^2$ and
  \Cref{cor:pet-ch6-riccati-comparison-estimate}(2). Interchanging $k$ and $K$
  fails already on flat $\mathbb{R}^2$, where the proposed upper bound becomes
  negative for $r>\pi/(2\sqrt K)$.
\end{theorem}
\begin{proof}
  \uses{cor:pet-ch6-riccati-comparison-estimate, rem:pet-ch6-jacobi-hessian-r}
  \sketch
  The radial Jacobi/Riccati equation gives \(\dot\rho\ge-K-\rho^2\) and \(\dot\rho\le-k-\rho^2\), with \(\rho=t^{-1}+O(t)\). Riccati comparison yields the Hessian bounds and the modified-distance inequalities.
\end{proof}

\begin{remark}[index-form proof of Rauch comparison]
  \label{rem:pet-ch6-index-form-remark}
  \lean{PetersenLib.rauchComparisonViaIndexForm,PetersenLib.rauchComparisonViaIndexForm_jacobiBoundary}
  \notready
  \source{petersen:page-0272}
  \uses{thm:pet-ch6-rauch-comparison}
  \dcref{ch6:6.4.4}
  A more traditional proof of the lower-bound half of
  \cref{thm:pet-ch6-rauch-comparison} uses the index form (exercise 6.7.25),
  and can also be adapted to the upper-bound case.
\end{remark}

\begin{example}[conjugate points in $S^n_K$]
  \label{ex:pet-ch6-conjugate-point-space-form}
  \lean{PetersenLib.conjugatePointSpaceFormExample, PetersenLib.isLeast_positive_zero_snFunction}
  \source{petersen:page-0273}
  \leanok
  \dcref{ch6:6.4.5}
  On $S^n_K$, $K>0$, polar coordinates around $p$ give
  $dr^2+\operatorname{sn}_K^2(r)\,d\sigma_{n-1}^2$; at distance $\pi/\sqrt K$ from $p$,
  every direction hits a conjugate point.

  The formalized content is the analytic core: $\operatorname{sn}_K$ solves the
  Jacobi equation $\ddot x+Kx=0$, and $\pi/\sqrt K$ is its \emph{first} positive
  zero (its full zero set being $\{n\pi/\sqrt K\}$), so the sphere directions of
  the polar form degenerate exactly there and nowhere earlier — which is what
  ``every direction hits a conjugate point at distance $\pi/\sqrt K$'' says. The
  companion fact $\operatorname{sn}_K(t)\ne0$ for $t>0$ when $K\le0$ is the
  analytic content of \cref{rem:pet-ch6-injectivity-radius-nonpositive}
  (no conjugate points in nonpositive curvature).
\end{example}

\begin{theorem}[conjugate-radius comparison]
  \label{thm:pet-ch6-conjugate-radius-bound}
  \lean{PetersenLib.conjugateRadiusLowerBound,PetersenLib.conjugateRadius_scalarEndpointUpperBound,PetersenLib.conjugateRadiusLowerBound_of_jacobiDomainCertificates}
  \notready
  \source{petersen:page-0273}
  \uses{thm:pet-ch6-rauch-comparison, ex:pet-ch6-conjugate-point-space-form}
  \dcref{ch6:6.4.6}
  If $(M,g)$ has $\sec\le K$, $K>0$, then
  $\exp_p:B(0,\pi/\sqrt K)\to M$ has no critical points.
\end{theorem}
\begin{proof}
  \uses{thm:pet-ch6-rauch-comparison, rem:pet-ch6-jacobi-hessian-r}
  \sketch
  Rauch comparison makes \(|J(t)|/\operatorname{sn}_K(t)\) nondecreasing with positive limit \(|\dot J(0)|\) at \(0\). Therefore no nontrivial \(J\) vanishes before \(\pi/\sqrt K\), so \(D\exp_p\) is nonsingular there.
\end{proof}

\begin{remark}[injectivity radius in nonpositive curvature]
  \label{rem:pet-ch6-injectivity-radius-nonpositive}
  \lean{PetersenLib.injectivityRadius_nonpositiveCurvature,PetersenLib.snFunction_noPositiveZero_nonpositiveCurvature,PetersenLib.injectivityRadius_nonpositiveCurvature_modelCore}
  \notready
  \source{petersen:page-0274}
  \uses{thm:pet-ch6-cartan-hadamard}
  If $\sec\le0$ there are no conjugate points, so
  $\operatorname{inj}(p)=\tfrac12\cdot(\text{length of shortest geodesic loop at }p)$.
  On a closed manifold with $\sec\le0$,
  $\operatorname{inj}(M)=\inf_p\operatorname{inj}(p)=\tfrac12\cdot(\text{length of shortest closed
  geodesic})$: the infimum is attained by compactness and continuity of
  $p\mapsto\operatorname{inj}(p)$; if $p$ realizes it via a geodesic loop $c:[0,1]\to M$,
  splitting $c$ at its midpoint $c(\tfrac12)$ gives
  $\operatorname{inj}(c(\tfrac12))\le\operatorname{inj}(p)$, forcing $c$ to also be a geodesic loop
  from $c(\tfrac12)$, hence smooth at $p$ and a closed geodesic.
\end{remark}

\begin{lemma}[Klingenberg's injectivity-radius lemma]
  \label{lem:pet-ch6-klingenberg-injectivity-estimate}
  \lean{PetersenLib.klingenbergInjectivityEstimate,PetersenLib.klingenbergInjectivityEstimate_of_domainCertificates}
  \notready
  \source{petersen:page-0274}
  \uses{thm:pet-ch6-conjugate-radius-bound, rem:pet-ch6-injectivity-radius-nonpositive}
  \dcref{ch6:6.4.7}
  Let $(M,g)$ be compact with $\sec\le K$, $K>0$. Then
  \[
    \operatorname{inj}(p)\ge\min\Big\{\frac{\pi}{\sqrt K},\ \tfrac12\cdot
    (\text{length of shortest geodesic loop at }p)\Big\},
  \]
  and $\operatorname{inj}(M)\ge\pi/\sqrt K$, or
  $\operatorname{inj}(M)=\tfrac12\cdot(\text{length of shortest closed geodesic})$.
\end{lemma}

\begin{theorem}[convexity-radius criterion]
  \label{thm:pet-ch6-convexity-radius-criterion}
  \lean{PetersenLib.convexityRadiusCriterion,PetersenLib.convexityRadiusCriterion_of_secondDerivative}
  \notready
  \source{petersen:page-0274}
  \uses{thm:pet-ch6-rauch-comparison}
  \dcref{ch6:6.4.8}
  Suppose $R>0$ satisfies (1) $R\le\tfrac12\operatorname{inj}(x)$ for all $x\in B(p,R)$,
  and (2) $R\le\tfrac12\cdot\pi/\sqrt K$, where
  $K=\sup\{\sec(\pi)\mid\pi\subset T_xM,\,x\in B(p,R)\}$. Then $r(x)=|xp|$ is
  convex on $B(p,R)$, and any two points of $B(p,R)$ are joined by a unique
  segment lying in $B(p,R)$.
\end{theorem}
\begin{proof}
  \uses{thm:pet-ch6-rauch-comparison}
  \sketch
  Injectivity gives unique segments and smooth \(r\), while Rauch gives \(\operatorname{Hess}r\ge0\). The set of endpoints whose segment from \(x\) stays in the ball is open; convexity of \(r\) makes it closed, hence all of the connected ball.
\end{proof}

\begin{definition}[convexity radius]
  \label{def:pet-ch6-convexity-radius}
  \lean{PetersenLib.convexityRadius}
  \source{petersen:page-0275}
  \uses{def:pet-ch6-convex-function, def:pet-ch5-segment, def:pet-ch5-metric-balls}
  \leanok
  The \emph{convexity radius} at $p$, $\operatorname{conv.rad}(p)$, is the largest
  $R>0$ such that $r(x)=|xp|$ is convex on $B(p,R)$ and any two points in $B(p,R)$
  are joined by a unique segment lying in $B(p,R)$. Globally
  $\operatorname{conv.rad}(M,g)=\inf_p\operatorname{conv.rad}(p)$.

  \emph{Formalization notes.} This is Petersen's own definition --- the two
  geometric properties themselves, not the sufficient hypotheses of
  \cref{thm:pet-ch6-convexity-radius-criterion} ($R\le\tfrac12\operatorname{inj}$,
  $R\le\tfrac12\pi/\sqrt K$). Encoding those hypotheses would define a smaller
  quantity, misname it, and needlessly carry a connection and $\sqrt K$
  junk-value handling; the $\ge\min\{\operatorname{inj}/2,\pi/(2\sqrt K)\}$
  estimate (the hypotheses-imply-conclusions direction) is the separate, still
  open \cref{cor:pet-ch6-convexity-radius-bound}. Both defining properties are
  metric, so \texttt{convexityRadius} takes \emph{no} connection argument and has
  the same shape as $\operatorname{inj}$. Two Lean-specific points: the value is an
  $\mathbb{R}_{\ge0}^\infty$ built as a $\sup$ over $\mathrm{ofReal}$-images
  (mirroring $\operatorname{inj}$), and ``unique segment'' is rendered as existence
  plus $\mathrm{EqOn}$ on the parameter interval $[0,|xy|]$, never as $\exists!$
  over $\mathbb{R}\to M$ --- the latter is false for every pair (a segment is
  unconstrained off its interval) and would silently make
  $\operatorname{conv.rad}\equiv0$. See \texttt{UniqueSegmentsInBall},
  \texttt{IsConvexityRadiusWitness}, \texttt{convexityRadius},
  \texttt{globalConvexityRadius} in \texttt{Ch06/ConvexityRadius.lean}.
\end{definition}

\begin{corollary}[convexity-radius lower bound]
  \label{cor:pet-ch6-convexity-radius-bound}
  \source{petersen:page-0275}
  \uses{def:pet-ch6-convexity-radius, thm:pet-ch6-convexity-radius-criterion, lem:pet-ch6-klingenberg-injectivity-estimate}
  \notready
  \notready
  By \cref{thm:pet-ch6-convexity-radius-criterion} and
  \cref{lem:pet-ch6-klingenberg-injectivity-estimate},
  $\operatorname{conv.rad}(M,g)\ge\min\{\operatorname{inj}(M,g)/2,\ \pi/(2\sqrt K)\}$ where
  $K=\sup\sec(M,g)$; in nonpositive curvature this simplifies to
  $\operatorname{conv.rad}(M,g)=\operatorname{inj}(M,g)/2$. (Split out from
  \cref{def:pet-ch6-convexity-radius} so that definition's Lean does not depend on
  the still-open \cref{thm:pet-ch6-convexity-radius-criterion} and
  \cref{lem:pet-ch6-klingenberg-injectivity-estimate}.)
\end{corollary}

\section{More on Positive Curvature}

\begin{theorem}[Klingenberg's even-dimensional injectivity estimate]
  \label{thm:pet-ch6-klingenberg-even-dim-injectivity}
  \lean{PetersenLib.klingenbergEvenDimInjectivity_of_data,PetersenLib.hasInjectivityLowerBound_iff_ofReal_le_globalInjectivityRadius}
  \notready
  \source{petersen:page-0276}
  \uses{lem:pet-ch6-klingenberg-injectivity-estimate, thm:pet-ch6-synge-theorem}
  \dcref{ch6:6.5.1}
  If $(M,g)$ is a compact orientable even-dimensional manifold with
  $0<\sec\le1$, then $\operatorname{inj}(M,g)\ge\pi$. If $M$ is not orientable,
  $\operatorname{inj}(M,g)\ge\pi/2$.
\end{theorem}
\begin{proof}
  \uses{lem:pet-ch6-klingenberg-injectivity-estimate, rem:pet-ch6-closed-parallel-field-construction}
  \sketch
  First assume that $M$ is orientable. The conjugate-radius estimate and
  $\sec\le1$ give conjugate radius at least $\pi$. If
  $\operatorname{inj}M<\pi$, Klingenberg's alternative therefore produces a
  closed geodesic $c$ of length $2\operatorname{inj}M<2\pi$. Even dimension and
  orientability give a closed parallel normal field along $c$ by
  \cref{rem:pet-ch6-closed-parallel-field-construction}; the Synge variation
  shortens $c$ to nearby closed curves $c_\varepsilon$. Join
  $c_\varepsilon(0)$ to a farthest point of $c_\varepsilon$ by a minimizing
  segment. At the farthest point this segment is perpendicular to
  $c_\varepsilon$, and a limit as $\varepsilon\to0$ is a minimizing segment
  from $c(0)$ to $c(\operatorname{inj}M)$ perpendicular to $c$ there. Below
  the conjugate radius the only minimizing segments between these endpoints
  are the two halves of $c$, neither perpendicular at their meeting point.
  This contradiction proves $\operatorname{inj}M\ge\pi$.

  If $M$ is nonorientable, apply the orientable result to its orientation
  double cover. The standard injectivity estimate for a double quotient gives
  $\operatorname{inj}M\ge\tfrac12\operatorname{inj}\widehat M\ge\pi/2$;
  equivalently, a shorter loop in $M$ would lift either to a loop of the same
  length or, after two turns, to a loop of length below $2\pi$ in
  $\widehat M$, contradicting its injectivity bound.
\end{proof}

\begin{remark}[odd-dimensional counterexamples]
  \label{rem:pet-ch6-klingenberg-odd-dim-counterexample}
  \lean{PetersenLib.klingenbergOddDimCounterexample,PetersenLib.bergerCounterexample_parameters_quarter,PetersenLib.klingenbergOddDimCounterexample_explicit,PetersenLib.klingenbergOddDimCounterexample_explicit_with_parameters}
  \notready
  \source{petersen:page-0276,petersen:page-0277}
  \uses{thm:pet-ch6-klingenberg-even-dim-injectivity}
  No analogue of \cref{thm:pet-ch6-klingenberg-even-dim-injectivity} holds in
  odd dimensions: the orientable quotients $S^3/\mathbb{Z}_k$ have a closed
  geodesic (image of the Hopf fiber) of length $2\pi/k\to0$; the Berger
  spheres $(S^3,g_\varepsilon)$ have Hopf fiber a closed geodesic of length
  $2\pi\varepsilon$ with curvatures in $[\varepsilon^2,4-3\varepsilon^2]$, so
  after rescaling the upper bound to $1$ the curvatures lie in
  $[\varepsilon^2/(4-3\varepsilon^2),1]$ and the Hopf fiber has length
  $2\pi\varepsilon\sqrt{4-3\varepsilon^2}<2\pi$ once $\varepsilon<1/\sqrt3$,
  with lower curvature bound then below $1/3$.
\end{remark}

\begin{definition}[$l$-connected pair]
  \label{def:pet-ch6-l-connected}
  \lean{PetersenLib.IsLConnected}
  \source{petersen:page-0277}
  $A\subset X$ is \emph{$l$-connected} if the relative homotopy groups
  $\pi_k(X,A)$ vanish for $k<l$; by Hurewicz's theorem the relative homology
  groups $H_k(X,A)$ then also vanish for $k\le l$, so from the long exact
  sequences of the pair, $\pi_k(A)\to\pi_k(X)$ and $H_k(A)\to H_k(X)$ are
  isomorphisms for $k<l$ and surjective for $k=l$.
\end{definition}

\begin{definition}[index of a critical point]
  \label{def:pet-ch6-critical-point-index}
  \lean{PetersenLib.criticalPointIndex}
  \source{petersen:page-0277}
  A critical point $p\in M$ of a smooth function $f:M\to\mathbb{R}$ has \emph{index}
  $m$ if $\operatorname{Hess}f$ is negative definite on an $m$-dimensional subspace of
  $T_pM$. If $m\ge1$, $p$ cannot be a local minimum of $f$, since $f$
  decreases along directions where the Hessian is negative definite.
\end{definition}

\begin{theorem}[Morse connectivity of sublevel sets]
  \label{thm:pet-ch6-morse-connectivity}
  \lean{PetersenLib.morseTheoreticConnectivity_of_data}
  \notready
  \source{petersen:page-0277}
  \uses{def:pet-ch6-l-connected, def:pet-ch6-critical-point-index}
  \dcref{ch6:6.5.2}
  Let $f:M\to\mathbb{R}$ be smooth and proper. If $b$ is not a critical value and
  every critical point in $f^{-1}([a,b])$ has index $\ge m$, then
  $f^{-1}((-\infty,a))\subset f^{-1}((-\infty,b))$ is $(m-1)$-connected.
\end{theorem}
\begin{proof}
  \sketch
  The Morse deformation removes critical points of index below \(m\): avoid each \(m\)-dimensional negative disk by transversality and push radially out. Every relative sphere of dimension \(<m\) therefore deforms below level \(a\), proving \((m-1)\)-connectivity.
\end{proof}

\begin{definition}[the path space $\Omega_{A,B}(M)$ and index of a geodesic]
  \label{def:pet-ch6-path-space-AB}
  \lean{PetersenLib.pathSpaceAB,PetersenLib.geodesicIndexAB}
  \source{petersen:page-0278}
  \uses{def:pet-ch6-critical-point-index}
  \leanok
  For $A,B\subset M$, $\Omega_{A,B}(M)=\{c:[0,1]\to M\mid c(0)\in A,\,c(1)\in
  B\}$. When $A,B$ are compact submanifolds, the energy functional
  $E:\Omega_{A,B}(M)\to[0,\infty)$ behaves as a proper smooth function, its
  critical points are geodesics meeting $A,B$ orthogonally at the endpoints,
  and its variational fields along such a geodesic are those tangent to $A,B$
  at the respective endpoints. The \emph{index} of such a geodesic is $\ge k$
  if there is a $k$-dimensional space of such variational fields on which the
  second variation of energy is negative (in the sense of
  \cref{def:pet-ch6-critical-point-index} applied to $E$ on
  $\Omega_{A,B}(M)$).
\end{definition}

\begin{theorem}[submanifold connectivity from geodesic index]
  \label{thm:pet-ch6-submanifold-connectivity}
  \lean{PetersenLib.submanifoldIndexConnectivity_of_data}
  \notready
  \source{petersen:page-0278}
  \uses{def:pet-ch6-path-space-AB, thm:pet-ch6-morse-connectivity}
  \dcref{ch6:6.5.3}
  Let $M$ be a complete Riemannian manifold and $A\subset M$ a compact
  submanifold. If every geodesic in $\Omega_{A,A}(M)$ perpendicular to $A$ at
  both endpoints has index $\ge k$, then $A\subset M$ is $k$-connected.
\end{theorem}
\begin{proof}
  \uses{thm:pet-ch6-morse-connectivity}
  \sketch
  Apply \cref{thm:pet-ch6-morse-connectivity} to the energy on \(\Omega_{A,A}\). Its critical points are geodesics perpendicular to \(A\), and the index bound gives \((k-1)\)-connectivity of constant paths; \(\pi_\ell(\Omega_{A,A},A)=\pi_{\ell+1}(M,A)\) gives \(k\)-connectivity of \(A\subset M\).
\end{proof}

\begin{theorem}[Berger's sphere theorem]
  \label{thm:pet-ch6-berger-sphere-theorem}
  \lean{PetersenLib.bergerSphereTheorem_of_data}
  \notready
  \source{petersen:page-0279}
  \uses{thm:pet-ch6-submanifold-connectivity, lem:pet-ch6-bonnet-synge-diameter, def:pet-ch6-l-connected}
  \dcref{ch6:6.5.4}
  Let $M$ be a closed $n$-manifold with $\sec\ge1$. If $\operatorname{inj}_p>\pi/2$ for
  some $p\in M$, then $M$ is $(n-1)$-connected, hence a homotopy sphere.
\end{theorem}
\begin{proof}
  \uses{thm:pet-ch6-submanifold-connectivity, lem:pet-ch6-bonnet-synge-diameter}
  \sketch
  At a point \(p\), loops have length \(>\pi\) and the Bonnet--Synge variation gives \(n-1\) negative directions. By \cref{thm:pet-ch6-submanifold-connectivity}, \(\{p\}\subset M\) is \((n-1)\)-connected; a degree-one map to \(S^n\), Hurewicz, and Whitehead give a homotopy sphere.
\end{proof}

\begin{theorem}[Klingenberg's pinching estimate]
  \label{thm:pet-ch6-klingenberg-pinching-injectivity}
  \lean{PetersenLib.klingenbergPinchingInjectivity_of_data}
  \notready
  \source{petersen:page-0279}
  \uses{thm:pet-ch6-klingenberg-even-dim-injectivity, lem:pet-ch6-klingenberg-injectivity-estimate}
  \dcref{ch6:6.5.5}
  A compact simply connected Riemannian $n$-manifold $(M,g)$ with
  $1\le\sec<4$ has $\operatorname{inj}>\pi/2$.
\end{theorem}
\begin{proof}
  \uses{thm:pet-ch6-klingenberg-even-dim-injectivity, lem:pet-ch6-klingenberg-injectivity-estimate}
  \sketch
  Rescale to \(1<\sec\le4\). In odd dimension, index bounds make every short loop homotopic through loops of length \(<\pi\); lifting this homotopy to the \(\pi/2\) exponential ball is open and closed, but the minimizing loop cannot lift closed. Hence \(\operatorname{inj}>\pi/2\); the even case is Klingenberg's estimate.
\end{proof}

\begin{corollary}[quarter-pinched sphere theorem]
  \label{cor:pet-ch6-rauch-berger-klingenberg-sphere}
  \lean{PetersenLib.rauchBergerKlingenbergSphereTheorem_of_data}
  \notready
  \source{petersen:page-0280}
  \uses{thm:pet-ch6-berger-sphere-theorem, thm:pet-ch6-klingenberg-pinching-injectivity}
  \dcref{ch6:6.5.6}
  Let $M$ be a closed simply connected $n$-manifold with $4>\sec\ge1$. Then
  $M$ is $(n-1)$-connected, hence a homotopy sphere.
\end{corollary}
\begin{proof}
  \uses{thm:pet-ch6-berger-sphere-theorem, thm:pet-ch6-klingenberg-pinching-injectivity}
  \sketch
  By \cref{thm:pet-ch6-klingenberg-pinching-injectivity}, \(\operatorname{inj}>\pi/2\); then \cref{thm:pet-ch6-berger-sphere-theorem} gives \((n-1)\)-connectivity and a homotopy sphere.
\end{proof}

\begin{corollary}[Ricci and injectivity imply simple connectivity]
  \label{cor:pet-ch6-ricci-simply-connected}
  \lean{PetersenLib.ricciDiameterInjectivitySimplyConnected_of_data}
  \notready
  \source{petersen:page-0280}
  \uses{thm:pet-ch6-submanifold-connectivity, thm:pet-ch6-myers-ricci-diameter}
  \dcref{ch6:6.5.7}
  If $M$ is a closed $n$-manifold with $\operatorname{Ric}\ge(n-1)$ and $\operatorname{inj}_p>\pi/2$
  for some $p\in M$, then $M$ is simply connected.
\end{corollary}
\begin{proof}
  \uses{thm:pet-ch6-submanifold-connectivity, thm:pet-ch6-myers-ricci-diameter}
  \sketch
  The Ricci bound and summed second variation give every nonconstant loop at \(p\) index at least one. Submanifold connectivity with \(A={p}\) yields simple connectivity.
\end{proof}

\begin{lemma}[the connectedness principle]
  \label{lem:pet-ch6-wilking-connectedness-principle}
  \lean{PetersenLib.wilkingConnectednessPrinciple_of_data}
  \notready
  \source{petersen:page-0280,petersen:page-0281}
  \uses{def:pet-ch6-l-connected, rem:pet-ch6-second-variation-special-cases}
  \dcref{ch6:6.5.8}
  Let $M^n$ be compact with $\sec>0$. (a) If $N^{n-k}\subset M^n$ is a closed
  totally geodesic codimension-$k$ submanifold, then $N\subset M$ is
  $(n-2k+1)$-connected. (b) If $N_1^{n-k_1},N_2^{n-k_2}\subset M$ are closed
  totally geodesic submanifolds with $k_1\le k_2$, $k_1+k_2\le n$, then
  $N_1\cap N_2$ is a nonempty totally geodesic submanifold and
  $N_1\cap N_2\to N_2$ is $(n-k_1-k_2)$-connected.
\end{lemma}
\begin{proof}
  \uses{thm:pet-ch6-submanifold-connectivity, rem:pet-ch6-second-variation-special-cases}
  \sketch
  For a geodesic between totally geodesic submanifolds, endpoint-tangent parallel fields have negative second variation. Dimension counting in \(\dot c^\perp\) proves the index bound in (a); applying \cref{thm:pet-ch6-submanifold-connectivity}, the same count forces intersection and gives the connectivity in (b) through path-space connectivity and retraction onto \(N_2\).
\end{proof}

\section{Further Study}

\begin{remark}[Frankel's theorem]
  \label{rem:pet-ch6-frankel-theorem}
  \lean{PetersenLib.frankelTheoremRemark}
  \source{petersen:page-0282}
  \uses{lem:pet-ch6-wilking-connectedness-principle}
  Part (b) of \cref{lem:pet-ch6-wilking-connectedness-principle}, in the
  special case that only nonemptiness of $N_1\cap N_2$ is retained, is
  classically known as \emph{Frankel's theorem}: two closed totally geodesic
  submanifolds $N_1,N_2$ of a manifold of positive curvature with
  $\dim N_1+\dim N_2\ge\dim M$ must intersect.
\end{remark}

\section{Exercises}

\begin{remark}[classification of even-dimensional spherical space forms]
  \label{rem:pet-ch6-ex-6-7-1}
  \lean{PetersenLib.exercise_6_7_1}
  \source{petersen:page-0282}
  \dcref{ch6:6.7.1}
  Show that in even dimensions the sphere and real projective space are the
  only closed manifolds with constant positive curvature.
\end{remark}

\begin{remark}[parallel transport on tangent cones]
  \label{rem:pet-ch6-ex-6-7-2}
  \lean{PetersenLib.exercise_6_7_2}
  \source{petersen:page-0282}
  \uses{def:pet-ch6-parallel-field}
  \dcref{ch6:6.7.2}
  Consider a rotationally symmetric metric $dr^2+\rho^2(r)d\theta^2$ and the
  cone $dr^2+(\rho(a)+\dot\rho(a)(r-a))^2d\theta^2$ tangent to it at the
  latitude $r=a$ (a genuine tangent cone when the surface is a surface of
  revolution). (1) Show that $\nabla_{\partial_\theta}$ agrees along $r=a$ on
  the two surfaces in the coordinates $(r,\theta)$, so parallel translation
  along $r=a$ is the same on both. (2) Unwrapping the cone to the flat plane,
  observe (by wrapping back up) that parallel translation along a latitude on
  a cone need not close up into a closed parallel field. (3) Show that the
  parallel field along $r=a$ closes up precisely when $\dot\rho(a)=0$.
\end{remark}

\begin{remark}[Fermi--Walker transport]
  \label{rem:pet-ch6-ex-6-7-3}
  \lean{PetersenLib.exercise_6_7_3}
  \source{petersen:page-0283}
  \uses{def:pet-ch6-vector-field-along-curve, def:pet-ch6-parallel-field}
  \dcref{ch6:6.7.3}
  Let $c:[a,b]\to M$ have nonvanishing speed and $T=\dot c/|\dot c|$. Call $V$
  along $c$ a \emph{Fermi--Walker field} if
  $\dot V=g(V,T)\dot T-g(V,\dot T)T=(T\wedge\dot T)(V)$. (1) Show that given
  $V(t_0)$ there is a unique Fermi--Walker field with that value. (2) Show
  $T$ itself is a Fermi--Walker field. (3) Show $g(V,W)$ is constant along
  $c$ for Fermi--Walker fields $V,W$. (4) Show that if $c$ is a geodesic,
  Fermi--Walker fields are parallel.
\end{remark}

\begin{remark}[coverings of constant-curvature manifolds]
  \label{rem:pet-ch6-ex-6-7-4}
  \lean{PetersenLib.exercise_6_7_4}
  \source{petersen:page-0283}
  \uses{thm:pet-ch6-conjugate-radius-bound}
  \dcref{ch6:6.7.4}
  Let $(M,g)$ be complete of constant curvature $k$ and $L:T_pM\to
  T_pS^n_k$ a linear isometry. If $k\le0$, show
  $\exp_p\circ L^{-1}\circ\exp_p^{-1}:S^n_k\to M$ is a Riemannian covering
  map. If $k>0$, show
  $\exp_p\circ L^{-1}\circ\exp_p^{-1}:S^n_k-\{-\tilde p\}\to M$ extends to a
  Riemannian covering map $S^n_k\to M$.
\end{remark}

\begin{remark}[flatness and two-parameter parallel transport]
  \label{rem:pet-ch6-ex-6-7-5}
  \lean{PetersenLib.exercise_6_7_5}
  \source{petersen:page-0283}
  \uses{def:pet-ch6-parallel-field}
  \dcref{ch6:6.7.5}
  Let $c(s,t):[0,1]^2\to(M,g)$ satisfy
  $R(\partial c/\partial s,\partial c/\partial t)=0$. Show that for each
  $v\in T_{c(0,0)}M$ there is a parallel field $V:[0,1]^2\to TM$ along $c$
  (i.e.\ $\partial V/\partial s=\partial V/\partial t=0$ everywhere) with
  $V(0,0)=v$.
\end{remark}

\begin{remark}[curvature symmetries from the commutator]
  \label{rem:pet-ch6-ex-6-7-6}
  \lean{PetersenLib.exercise_6_7_6}
  \source{petersen:page-0283}
  \uses{lem:pet-ch6-third-partial-commutator}
  \dcref{ch6:6.7.6}
  Using
  $R(\partial c/\partial s,\partial c/\partial t)\partial c/\partial u
  =\partial^3c/\partial s\,\partial t\,\partial u
  -\partial^3c/\partial t\,\partial s\,\partial u$, show that the two
  skew-symmetries and the first Bianchi identity for the curvature tensor
  (proposition 3.1.1 of the source) hold.
\end{remark}

\begin{remark}[Killing fields along geodesics]
  \label{rem:pet-ch6-ex-6-7-7}
  \lean{PetersenLib.exercise_6_7_7}
  \source{petersen:page-0284}
  \uses{def:pet-ch6-jacobi-field}
  \dcref{ch6:6.7.7}
  Let $c$ be a geodesic and $X$ a Killing field. Show $X|_c$ is a Jacobi
  field.
\end{remark}

\begin{remark}[conjugate points and Jacobi fields]
  \label{rem:pet-ch6-ex-6-7-8}
  \lean{PetersenLib.exercise_6_7_8}
  \source{petersen:page-0284}
  \uses{def:pet-ch6-jacobi-field, rem:pet-ch6-jacobi-dexp-relation}
  \dcref{ch6:6.7.8}
  Let $c:[0,1]\to M$ be a geodesic. Show $\exp_{c(0)}$ has a critical point at
  $\dot c(0)$ iff there is a nontrivial Jacobi field $J$ along $c$ with
  $J(0)=0$, $\dot J(0)\perp\dot c(0)$, $J(1)=0$.
\end{remark}

\begin{remark}[Hessian of squared distance]
  \label{rem:pet-ch6-ex-6-7-9}
  \lean{PetersenLib.exercise_6_7_9}
  \source{petersen:page-0284}
  \uses{def:pet-ch6-jacobi-field, lem:pet-ch6-hessian-lower-bound}
  \dcref{ch6:6.7.9}
  Fix $p$, $v\in\operatorname{seg}_p^0$, $c(t)=\exp_p(tv)$, and the geodesic variation
  $\bar c(s,t)=\exp_p(t(v+sw))$ with Jacobi field $J$. Show
  $\nabla f_0|_{c(1)}=\dot c(1)$ and
  $\operatorname{Hess}f_0(J(1),J(1))=g(\dot J(1),\dot J(1))$ for $f_0=\tfrac12r^2$; use
  this to prove \cref{lem:pet-ch6-hessian-lower-bound} without first
  estimating $\operatorname{Hess}r$.
\end{remark}

\begin{remark}[Wronskian identities for Jacobi fields]
  \label{rem:pet-ch6-ex-6-7-10}
  \lean{PetersenLib.exercise_6_7_10}
  \source{petersen:page-0284}
  \uses{def:pet-ch6-jacobi-field}
  \leanok
  \dcref{ch6:6.7.10}
  Let $c$ be a geodesic and $J_1,J_2$ Jacobi fields along $c$. (1) Show
  $g(\dot J_1,J_2)-g(J_1,\dot J_2)$ is constant. (2) Show
  $g(J_1(t),\dot c(t))=g(J_1(0),\dot c(0))+g(\dot J_1(0),\dot c(0))t$.
\end{remark}

\begin{remark}[a Riccati inequality for Jacobi fields]
  \label{rem:pet-ch6-ex-6-7-11}
  \lean{PetersenLib.exercise_6_7_11}
  \source{petersen:page-0284}
  \uses{def:pet-ch6-jacobi-field}
  \dcref{ch6:6.7.11}
  Let $J$ be a nontrivial Jacobi field along a unit speed geodesic $c$ with
  $J(0)=0$, $\dot J(0)\perp\dot c(0)$, and suppose $\sec\le K$. (1) With
  $\lambda=|J|^2/g(J,\dot J)$, show $\dot\lambda\le1+K\lambda^2$, $\lambda(0)=0$
  while $\lambda$ is defined. (2) Show that if $J(b)=0$ for some $b>0$, then
  $g(J(t),\dot J(t))=0$ for some $t\in(0,b)$; give an explicit example.
\end{remark}

\begin{remark}[self-adjoint Jacobi spaces]
  \label{rem:pet-ch6-ex-6-7-12}
  \lean{PetersenLib.exercise_6_7_12}
  \source{petersen:page-0284}
  \uses{def:pet-ch6-jacobi-field}
  \dcref{ch6:6.7.12}
  Let $\mathfrak J$ be a self-adjoint space of Jacobi fields along a geodesic
  $c$ ($g(\dot J_1,J_2)=g(J_1,\dot J_2)$ for $J_1,J_2\in\mathfrak J$), and
  $\tilde{\mathfrak J}(t)=\{J(t):J\in\mathfrak J,\,J(t)=0\}
  +\{J(t):J\in\mathfrak J\}\subset T_{c(t)}M$ — read as the ambient span,
  with the two summands orthogonal. (1) Show the two summands are orthogonal.
  (2) Show $\{J\in\mathfrak J:J(t)=0\}$ is naturally isomorphic to the first
  summand. (3) Show $\dim\tilde{\mathfrak J}=\dim\tilde{\mathfrak J}(t)$ for
  all $t$.
\end{remark}

\begin{remark}[point homogeneity]
  \label{rem:pet-ch6-ex-6-7-13}
  \lean{PetersenLib.exercise_6_7_13}
  \source{petersen:page-0285}
  \uses{def:pet-ch6-vector-field-along-curve}
  \dcref{ch6:6.7.13}
  Call $(M,g)$ \emph{$k$-point homogeneous} if for all $(p_1,\dots,p_k)$,
  $(q_1,\dots,q_k)$ with $|p_ip_j|=|q_iq_j|$ there is an isometry $F$ with
  $F(p_i)=q_i$ ($k=1$: homogeneous). (1) Show a homogeneous space has
  constant scalar curvature. (2) Show $k$-point homogeneous
  ($k>1$) implies $(k-1)$-point homogeneous. (3) Show two-point homogeneous
  implies Einstein. (4) Show three-point homogeneous implies constant
  curvature. (5) Show $\mathbb{RP}^2$ is not three-point homogeneous by
  exhibiting two noncongruent equilateral triangles of side length $\pi/3$.
\end{remark}

\begin{remark}[equidistant curves in space forms]
  \label{rem:pet-ch6-ex-6-7-14}
  \lean{PetersenLib.exercise_6_7_14}
  \source{petersen:page-0285}
  \dcref{ch6:6.7.14}
  Starting with a geodesic on a two-dimensional space form, discuss how the
  equidistant curves change as they move away from the geodesic.
\end{remark}

\begin{remark}[two-sided metric comparison]
  \label{rem:pet-ch6-ex-6-7-15}
  \lean{PetersenLib.exercise_6_7_15}
  \source{petersen:page-0285}
  \uses{def:pet-ch6-jacobi-field, thm:pet-ch6-rauch-comparison}
  \dcref{ch6:6.7.15}
  Let $r(x)=|xp|$ with $-K\le\sec\le K$, $g=dr^2+g_r$ on
  $(0,R)\times S^{n-1}$, $2R<\operatorname{inj}_p$. (1) Show
  $\operatorname{sn}_K^2(r)\,ds_{n-1}^2\le g_r\le\operatorname{sn}_{-K}^2(r)\,ds_{n-1}^2$. (2) Show there
  is a universal constant $C$ with
  $|\operatorname{Hess}\tfrac12r^2-g|\le CK^2R^2$ for $R<\pi/(2\sqrt K)$.
\end{remark}

\begin{remark}[shortest loops in a homotopy class]
  \label{rem:pet-ch6-ex-6-7-16}
  \lean{PetersenLib.exercise_6_7_16}
  \source{petersen:page-0285}
  \dcref{ch6:6.7.16}
  Let $(M,g)$ be complete. Show every element of $\pi_1(M,p)$ contains a
  shortest loop at $p$, and that this shortest loop is a geodesic loop.
\end{remark}

\begin{remark}[index of a shortest loop]
  \label{rem:pet-ch6-ex-6-7-17}
  \lean{PetersenLib.exercise_6_7_17}
  \source{petersen:page-0285}
  \uses{def:pet-ch6-path-space-AB}
  \dcref{ch6:6.7.17}
  Let $(M,g)$ be complete with $\operatorname{inj}_p<R$ where
  $\exp_p:B(0,R)\to B(p,R)$ is nonsingular. Show the geodesic loop $c$ at $p$
  realizing the injectivity radius has index $0$.
\end{remark}

\begin{remark}[Frankel intersection argument]
  \label{rem:pet-ch6-ex-6-7-18}
  \lean{PetersenLib.exercise_6_7_18}
  \source{petersen:page-0286}
  \uses{rem:pet-ch6-frankel-theorem}
  \dcref{ch6:6.7.18}
  Let $M$ be $n$-dimensional with $\sec>0$ and $A,B$ closed totally geodesic
  submanifolds. Show directly (via a shortest connecting segment and the
  second variation formula, without invoking
  \cref{lem:pet-ch6-wilking-connectedness-principle}) that $A\cap B\ne
  \emptyset$ if $\dim A+\dim B\ge n$.
\end{remark}

\begin{remark}[path-space proof of connectedness]
  \label{rem:pet-ch6-ex-6-7-19}
  \lean{PetersenLib.exercise_6_7_19}
  \source{petersen:page-0286}
  \uses{def:pet-ch6-path-space-AB, lem:pet-ch6-wilking-connectedness-principle}
  \dcref{ch6:6.7.19}
  Let $M$ be complete, $A\subset M$ compact, without invoking
  \cref{lem:pet-ch6-wilking-connectedness-principle}. (1) Show non-stationary
  curves in $\Omega_{A,A}(M)$ deform to shorter curves in $\Omega_{A,A}(M)$.
  (2) Show stationary curves are geodesics perpendicular to $A$ at the
  endpoints. (3) If $M$ has $\sec>0$, $A$ is totally geodesic, and $2\dim
  A\ge\dim M$, show all stationary curves deform to shorter curves in
  $\Omega_{A,A}(M)$. (4) (Wilking) Conclude that $\pi_1(M,A)$ is trivial.
\end{remark}

\begin{remark}[solvable subgroups in negative curvature]
  \label{rem:pet-ch6-ex-6-7-20}
  \lean{PetersenLib.exercise_6_7_20}
  \source{petersen:page-0286}
  \uses{thm:pet-ch6-preissmann, def:pet-ch6-axis-displacement}
  \dcref{ch6:6.7.20}
  Generalize \cref{thm:pet-ch6-preissmann} to show any solvable subgroup of
  $\pi_1$ of a compact negatively curved manifold is cyclic. Hint: use
  torsion-freeness and, by contradiction and solvability, a subgroup
  generated by deck transformations $F,G$ with $F\circ G=G^k\circ F$,
  $k\ne0$; show that if $c$ is an axis for $G$, $F\circ c$ is an axis for
  $G^k$, and use uniqueness of axes for $G^k$ for a contradiction.
\end{remark}

\begin{remark}[orientation of fixed-point-free isometries]
  \label{rem:pet-ch6-ex-6-7-21}
  \lean{PetersenLib.exercise_6_7_21}
  \source{petersen:page-0286}
  \uses{rem:pet-ch6-closed-parallel-field-construction}
  \dcref{ch6:6.7.21}
  Let $(M,g)$ be compact with $\sec>0$ and $F:M\to M$ an isometry of finite
  order without fixed points. Show that if $\dim M$ is even, $F$ is
  orientation reversing; if odd, orientation preserving. (Weinstein: this
  holds even without the finite-order hypothesis.)
\end{remark}

\begin{remark}[diameter of nontrivial spherical quotients]
  \label{rem:pet-ch6-ex-6-7-22}
  \lean{PetersenLib.exercise_6_7_22}
  \source{petersen:page-0286}
  \uses{thm:pet-ch6-cartan-fixed-point}
  \dcref{ch6:6.7.22}
  Use an analogue of \cref{thm:pet-ch6-cartan-fixed-point} to show a closed
  manifold of constant curvature $1$ is either the standard sphere or has
  diameter $\le\pi/2$; generalize to show a closed manifold with $\sec\ge1$
  is either simply connected or has diameter $\le\pi/2$.
\end{remark}

\begin{remark}[distance coordinates]
  \label{rem:pet-ch6-ex-6-7-23}
  \lean{PetersenLib.exercise_6_7_23}
  \source{petersen:page-0286,petersen:page-0287}
  \uses{thm:pet-ch6-rauch-comparison}
  \dcref{ch6:6.7.23}
  Let $(M,g)$ be complete $n$-dimensional with $|\sec|<K$, and fix $p$,
  $\epsilon>0$, points $p_i$ with $r^i(x)=|xp_i|$ smooth on $B(p,\epsilon)$,
  $g(\nabla r^i,\nabla r^j)|_p=\delta^{ij}$, $|pp_i|\ge2\epsilon$. Let
  $g^{ij}=g(\nabla r^i,\nabla r^j)$. (1) Show there is $C=C(n,K,\epsilon)$
  with $|dg^{ij}|\le C$. (2) Show $C$ can be chosen so that
  $C(n,\lambda^{-2}K,\lambda\epsilon)\to0$ as $\lambda\to\infty$. (3) Show
  there is $\delta=\delta(n,C)>0$ with $g^{ij}$ invertible on $B(p,\delta)$
  and its inverse $g_{ij}$ satisfying $|g_{ij}-\delta_{ij}|\le\tfrac12$,
  $|dg_{ij}|\le C'(n,K,\epsilon)$. (4) Show
  $(r^1(x)-r^1(p),\dots,r^n(x)-r^n(p))$ is a coordinate system on
  $B(p,\delta)$ whose image contains $B(0,\delta/2)$.
\end{remark}

\begin{remark}[the index form]
  \label{rem:pet-ch6-ex-6-7-24}
  \lean{PetersenLib.exercise_6_7_24,PetersenLib.indexForm}
  \source{petersen:page-0287,petersen:page-0288}
  \uses{rem:pet-ch6-second-variation-special-cases, def:pet-ch6-jacobi-field}
  \dcref{ch6:6.7.24}
  For $V,W$ along a geodesic $c:[0,1]\to(M,g)$, the \emph{index form} is
  \[
    I_0^1(V,W)=I(V,W)=\int_0^1\big(g(\dot V,\dot W)-g(R(V,\dot c)\dot c,W)\big)\,dt,
  \]
  agreeing with the second variation of energy when $V,W$ come from a proper
  variation. Assume $c$ locally minimizes energy, so $I(V,V)\ge0$ for proper
  variational fields. (1) If $I(V,V)=0$ for such $V$, show $V$ is a Jacobi
  field (expand $I(V+\epsilon W,V+\epsilon W)\ge0$ in $\epsilon$ for all
  proper $W$). (2) If $V(0)=J(0)$, $V(1)=J(1)$ and $J$ is a Jacobi field,
  show $I(V,J)=I(J,J)$. (3) \emph{(Index lemma)} If in addition no nontrivial
  Jacobi field vanishes at both endpoints of $c$, show $I(V,V)\ge I(J,J)$
  with equality iff $V=J$ (via $0\le I(V-J,V-J)=I(V,V)-I(J,J)$). (4) If a
  nontrivial Jacobi field $J$ vanishes at $0$ and $1$, show
  $c:[0,1+\epsilon]\to M$ is not locally minimizing for $\epsilon>0$, by
  patching $J$ on $[0,1-\epsilon]$ to a Jacobi field $K$ on $[1-\epsilon,
  1+\epsilon]$ with $K(1-\epsilon)=J(1-\epsilon)$, $K(1+\epsilon)=0$, and
  comparing $I_0^1(J,J)=0$ against $I_0^{1-\epsilon}(J,J)+I_{1-\epsilon}^{1+
  \epsilon}(K,K)$, which is strictly smaller since $K\ne J$ on
  $[1-\epsilon,1]$.
\end{remark}

\begin{remark}[index comparison]
  \label{rem:pet-ch6-ex-6-7-25}
  \lean{PetersenLib.exercise_6_7_25}
  \source{petersen:page-0288}
  \uses{rem:pet-ch6-ex-6-7-24, thm:pet-ch6-rauch-comparison}
  \dcref{ch6:6.7.25}
  Let $J$ be a nontrivial Jacobi field along a unit speed geodesic $c$ with
  $J(0)=0$, $\dot J(0)\perp\dot c(0)$, on a manifold with $\sec\ge k$, and
  suppose no Jacobi field on $c|_{[0,b]}$ vanishes at both endpoints. (1)
  Show $I_0^b(J,J)=g(J(b),\dot J(b))$. (2) With $V(t)=\frac{\operatorname{sn}_k(t)}
  {\operatorname{sn}_k(b)}E(t)$ for a parallel field $E$ with $E(b)=J(b)$, show
  $I_0^b(V,V)\le\frac{\operatorname{sn}_k'(b)}{\operatorname{sn}_k(b)}|J(b)|^2$. (3) Conclude
  $g(J(b),\dot J(b))\le\frac{\operatorname{sn}_k'(b)}{\operatorname{sn}_k(b)}|J(b)|^2$, and use this to
  prove the lower-curvature-bound half of
  \cref{thm:pet-ch6-rauch-comparison}.
\end{remark}

\begin{remark}[small subgroups of the isometry group]
  \label{rem:pet-ch6-ex-6-7-26}
  \lean{PetersenLib.exercise_6_7_26}
  \source{petersen:page-0288,petersen:page-0289}
  \uses{thm:pet-ch6-cartan-fixed-point, def:pet-ch6-convexity-radius}
  \dcref{ch6:6.7.26}
  Let $\mathbf G\subset\operatorname{Iso}(M,g)$, with $\operatorname{Iso}(M,g)$ given the
  compact-open topology. (1) If $M$ is complete, simply connected,
  nonpositively curved, and $\mathbf G$ is compact, show $\mathbf G$ has a
  fixed point (imitate \cref{thm:pet-ch6-cartan-fixed-point}). (2) Call
  $\mathbf G$ \emph{$(p,\epsilon)$-small} if $\mathbf Gp\subset\bar
  B(p,\epsilon)$; show that for sufficiently small $\epsilon(p)$ the closure
  $\bar{\mathbf G}$ of a $(p,\epsilon)$-small group is compact and
  $(p,\epsilon)$-small (without assuming $M$ complete). (3) Show a
  $(p,\epsilon)$-small $\mathbf G$ is $(q,2|pq|+\epsilon)$-small. (4) If
  $\epsilon$ is much smaller than the convexity radius on $\bar
  B(p,4\epsilon)$, show $\mathbf G$ has a fixed point. (5) Show that for
  every $p$ there is $\epsilon>0$ such that no subgroup of $\operatorname{Iso}(M,g)$ can
  be $(p,\epsilon)$-small (by making $\mathbf G$ act freely on a suitable
  subset of a product $M\times\cdots\times M$). (6) Conclude
  $\operatorname{Iso}(M,g)$ has no small subgroups (no neighborhood of the identity
  contains a nontrivial subgroup).
\end{remark}

\begin{remark}[a submersion metric on the tangent bundle]
  \label{rem:pet-ch6-ex-6-7-27}
  \lean{PetersenLib.exercise_6_7_27}
  \source{petersen:page-0289}
  \uses{def:pet-ch6-parallel-field}
  \dcref{ch6:6.7.27}
  Construct a Riemannian metric on $TM$ for a Riemannian manifold $(M,g)$
  such that $\pi:TM\to M$ is a Riemannian submersion restricting to the given
  Euclidean metric on each fiber, by declaring parallel fields along curves
  in $M$ to give horizontal lifts.
\end{remark}

\begin{remark}[a submersion metric on the frame bundle]
  \label{rem:pet-ch6-ex-6-7-28}
  \lean{PetersenLib.exercise_6_7_28}
  \source{petersen:page-0289}
  \uses{def:pet-ch6-parallel-field}
  \dcref{ch6:6.7.28}
  For a Riemannian manifold $(M,g)$ let $FM\to M$ be the orthonormal frame
  bundle. Find a Riemannian metric on $FM$ making $\pi:FM\to M$ a Riemannian
  submersion with fibers isometric to $O(n)$, by declaring parallel
  orthonormal frames along curves in $M$ to give horizontal lifts.
\end{remark}
